\chapter{多 GPU 编程}
\label{chap:multi-gpu-programming}
\chaptercontents

\setcounter{footnote}{0}

\noindent\textit{并介绍 MPI、NCCL 和 NVSHMEM。Jiri Kraus 和 Isaac Gelado 特别参与撰写。}

到目前为止，我们一直关注如何为一个主机和一个设备组成的异构计算系统编程。然而，许多
应用需要汇集多个主机和设备的计算能力。如今，许多集群和数据中心的每个节点都包含一个或
多个主机，并连接一个或多个设备。本章考察三种面向多 GPU 集群的编程模型：消息传递接口
（Message Passing Interface，MPI）\cite{ch23-mpi}、NVIDIA 集体通信库（NVIDIA
Collective Communications Library，NCCL）\cite{ch23-nccl}，以及 NVIDIA OpenSHMEM
库（NVIDIA OpenSHMEM Library，NVSHMEM）\cite{ch23-nvshmem}。

我们只介绍每种编程模型中，将异构应用扩展到集群内多个 GPU 和节点时必须理解的关键概念。
具体来说，我们以跨多个 GPU 节点扩展 CUDA 模板计算为背景，重点讨论域分区、点对点通信、
集体通信，以及计算与通信的重叠执行。本章给出的实现旨在说明每种编程模型背后的核心原则；
在各自模型中，它们都还有进一步优化的空间。掌握本章概念以后，建议读者查阅三种编程模型
各自的教程和应用案例，以便更全面地了解它们的全部能力。

\begin{figure}[htbp]
  \centering
  \begin{lstlisting}[language=C++,basicstyle=\ttfamily\scriptsize,firstnumber=1]
__global__ void jacobi_kernel(float* out, float* in, float* l2norm,
                              unsigned int ny, unsigned int nx) {
    int y = blockIdx.y*blockDim.y + threadIdx.y + 1;
    int x = blockIdx.x*blockDim.x + threadIdx.x + 1;
    // Compute stencil and residue
    float residue = 0.0f;
    if(y < ny - 1 && x < nx - 1) {
        float val = 0.25*(in[y*nx + (x + 1)]
                        + in[y*nx + (x - 1)]
                        + in[(y + 1)*nx + x]
                        + in[(y - 1)*nx + x]);
        out[y*nx + x] = val;
        residue = val - in[y*nx + x];
    }
    // Reduce L2 norm
    float blockL2norm = blockReduce(residue*residue);
    if (threadIdx.y == 0 && threadIdx.x == 0) {
        atomicAdd(l2norm, blockL2norm);
    }
}
  \end{lstlisting}
  \caption{二维 Jacobi 迭代法的修改版模板内核。依据原书图 23.1 逐字符重排。}
  \label{fig:23-1}
\end{figure}

\section{作为贯穿示例的模板计算}
\label{sec:multi-gpu-running-stencil}

本章始终以第 8 章介绍的模板计算作为贯穿示例。具体而言，我们采用 Jacobi 迭代法：在每次
迭代或每个时间步中，结构化网格\footnote{本章用“网格”表示第 \ref{sec:stencil-background}
节所介绍的、建模系统中由数据值组成的结构化网格。为免混淆，GPU 线程组成的网格称为
“线程网格”。}中一个点的值，等于上一时间步中其邻点值的加权和。模板内核会执行多次迭代，
直到网格点值不再发生显著变化。为判断这些值是否已经停止变化，我们计算残差值的 $L_2$
范数；每个网格点的残差，就是该点旧值与新值之差。

图 \ref{fig:23-1} 的模板内核相对于第 8 章的版本有两处修改\cite{ch23-multigpu-models}。
第一处修改是计算二维模板而不是三维模板，以便简化插图和代码示例。第二处修改是让模板内核
同时计算残差及其 $L_2$ 范数。图中，每个线程首先确定自己负责计算的网格点的 $x$、$y$
坐标（第 3--4 行）。边界检查保证只有负责界内网格点的线程保持活动（第 7 行）。注意，
$x=0$、$x=nx-1$、$y=0$ 和 $y=ny-1$ 处的点会被跳过，因为它们是边界点，或者是后文
将会看到的 halo 点。线程把相邻网格点的旧值取平均，得到自己负责网格点的新值（第 8--12
行）。到此为止的代码应当与第 8 章的内核十分相似。

计算出新值以后，每个线程用新值减去旧值，得到该网格点的残差（第 13 行）。计算 $L_2$
范数时，需要把线程网格中所有线程的残差平方加在一起。代码使用第 10 章讨论过的块范围归约，
求出整个块的残差平方和（第 16 行）；随后由块中的一个线程，以原子方式把块级平方和累加到
线程网格范围的 $L_2$ 范数平方和中（第 17--19 行）。

图 \ref{fig:23-1} 的内核还可以继续优化。例如，可以按照第 8 章的方法使用共享内存分块和
线程粗化，减少全局内存访问；还可以采用第 10 章介绍的多种归约优化，或调用 CUB
库\cite{ch23-cub}的块范围归约函数。不过，这些优化不是本章重点；本章关注如何把模板计算
分布到多个 GPU 上。

\begin{figure}[htbp]
  \centering
  \begin{tikzpicture}[x=0.75cm,y=0.58cm,>=Stealth]
    \foreach \y/\c/\g in {6.0/red!68/0,4.2/yellow!75/1,2.4/green!60/2,0.6/blue!55/3} {
      \fill[\c] (0,\y) rectangle (7,\y+1.5);
      \draw[thick] (0,\y) rectangle (7,\y+1.5);
      \draw[dashed] (0,\y+0.18) -- (7,\y+0.18);
      \draw[dashed] (0,\y+1.32) -- (7,\y+1.32);
      \node[left,font=\small] at (0,\y+0.75) {GPU \g};
    }
    \draw[<->] (0,8.1) -- (7,8.1) node[midway,above] {$nx$};
    \draw[<->] (7.35,4.2) -- (7.35,5.7) node[midway,right] {$ny$};
  \end{tikzpicture}
  \caption{跨 4 个 GPU 划分建模系统的二维网格数组。依据原书图 23.2 重绘；实线分隔各
  GPU 负责的域分区，虚线标出 halo 行。}
  \label{fig:23-2}
\end{figure}

图 \ref{fig:23-2} 展示了建模系统，以及如何把它划分到多个 GPU 上。假定系统由结构化
网格建模，各方向上相邻网格点的间距分别保持不变。根据这一假设，可以像第 8 章那样用二维
数组表示系统，每个元素保存一个网格点的状态；每个维度上相邻元素的物理距离可以用一个网格
间距变量表示。我们还假定二维数组在内存空间中采用行主序布局，$x$ 是最低维度，随后是
$y$。也就是说，所有 $y=0$ 的元素按照 $x$ 坐标放入连续内存位置，接着存放所有 $y=1$
的元素，依此类推。

使用多个 GPU 时，通常把输入数据划分成若干称为\textbf{域分区（domain partition）}的
部分，再把每个分区分配给集群中的不同 GPU。采用行主序布局时，沿 $y$ 维在 GPU 之间划分
系统很自然，因为每个 GPU 都会得到一组连续的行。图 \ref{fig:23-2} 用不同颜色表示各 GPU
负责计算的网格点，并用实线分隔各分区。这样划分以后，每个 GPU 的分区在行主序内存布局中
都是连续的。

回忆一下，计算每个网格点的新值时，需要相邻点的旧值。因此，一个 GPU 要计算其分区边缘
网格点的新值，就必须取得相邻分区边缘网格点的旧值。这些网格点称为\textbf{halo 点
（halo point）}，在图 \ref{fig:23-2} 中以虚线标出。例如，GPU 1 要计算自己分区内网格点
的新值，除了需要这些点的旧值，还需要 GPU 0 分区最底一行和 GPU 2 分区最顶一行的旧值。
GPU 1 所需的输入行跨度用 $ny$ 表示；其中第一行和最后一行是只作输入的 halo 点，中间的
$ny-2$ 行既是输入，也是输出。

\begin{figure}[htbp]
  \centering
  \begin{tikzpicture}[x=0.54cm,y=0.48cm,>=Stealth]
    \node[font=\small\bfseries] at (3.5,7.9) {(a) 计算前可用的旧值};
    \node[font=\small\bfseries] at (12.5,7.9) {(b) 计算后可用的新值};
    \foreach \base/\panel in {0/0,9/1} {
      \foreach \y/\c/\g in {5.8/red!68/0,3.3/yellow!75/1,0.8/green!60/2} {
        \fill[\c] (\base,\y) rectangle +(7,1.7);
        \draw[thick] (\base,\y) rectangle +(7,1.7);
        \draw[dashed] (\base,\y+0.17) -- +(7,0);
        \draw[dashed] (\base,\y+1.53) -- +(7,0);
        \node[left,font=\scriptsize] at (\base,\y+0.85) {GPU \g};
      }
    }
    \foreach \y in {5.8,3.3} {
      \draw[<->,orange!80!black] (16.2,\y-0.03) .. controls (16.8,\y-0.35) and
        (16.8,\y-0.65) .. (16.2,\y-0.87);
      \node[right,font=\scriptsize,align=left] at (16.55,\y-0.45) {halo\\交换};
    }
    \node[font=\scriptsize,align=center] at (3.5,-0.15) {各 GPU 已取得相邻分区的 halo 行};
    \node[font=\scriptsize,align=center] at (12.5,-0.15) {本地 halo 行已过期，必须重新交换};
  \end{tikzpicture}
  \caption{各 GPU 在计算前后可用的值。依据原书图 23.3 重绘。}
  \label{fig:23-3}
\end{figure}

图 \ref{fig:23-3} 进一步说明每个 GPU 在计算其分区网格点之前和之后分别拥有哪些值。计算
之前，GPU 需要图 \ref{fig:23-3}(a) 所示的 halo 行值。计算结束后，其 halo 行值会像图
\ref{fig:23-3}(b) 那样变得过期，因为这些行的新值由负责相应行的其他 GPU 计算。因此，
进入下一次迭代之前，GPU 必须从计算相邻分区的 GPU 取得更新后的 halo 行。例如，GPU 1
需要从 GPU 0 复制顶部 halo 行，从 GPU 2 复制底部 halo 行；与此同时，GPU 1 还要把自己
的顶部边界行发送给 GPU 0，把底部边界行发送给 GPU 2。相邻 GPU 之间这样发送和接收 halo
行，称为\textbf{halo 交换（halo exchange）}。

这里还能看到，在行主序布局下沿 $y$ 维进行一维域分解的另一项好处：GPU 之间需要复制的
halo 行在内存中连续，从而简化 halo 交换期间的复制操作。这种一维域分解也有缺点：根据域的
形状和 GPU 数量，GPU 之间可能需要传输大量数据。例如，如果域很宽，每个 GPU 都必须传输
完整的一行，代价可能很高。另一种选择是采用二维域分解，让每个 GPU 负责域中的一个二维
瓦片。二维域分解可以降低表面积体积比，从而减少需要传输的数据量；但它会要求每个 GPU 与
四个相邻 GPU 通信，而不再只是两个，并且部分通信数据在内存中并不连续。为保持本章示例
简单，下面只讨论所述的一维域分解。

总之，多 GPU Jacobi 迭代法的每次迭代包含两个阶段：一是计算阶段，求出网格点的新值；
二是通信阶段，在负责相邻分区的 GPU 之间交换新的 halo 行值。由于 $L_2$ 范数必须覆盖所有
分区内的全部网格点，通信阶段还包含一次归约操作，用来归约各 GPU 计算出的线程网格范围
$L_2$ 范数平方和。MPI、NCCL 和 NVSHMEM 分别提供了自己的 GPU 间数据交换机制；本章后续
将依次说明如何用这三种编程模型完成这些操作。

\section{使用 MPI 的多 GPU 模板计算}
\label{sec:multi-gpu-mpi}

并行编程模型主要分为两类：共享内存模型和分布式内存模型。在共享内存模型中，不同线程或
进程可以访问同一个内存空间，因而数据共享较为简单。例如，CUDA 内核采用共享内存模型，
不同 CUDA 线程可以用同一个指针访问同一个全局内存变量。在分布式内存模型中，不同线程或
进程访问各自独立的内存空间；此时必须通过消息传递等方式显式共享数据。

当今计算集群中占主导地位的编程接口是 MPI\cite{ch23-mpi}。MPI 是一组 API 函数，用于在
计算集群中运行的进程之间通信。它假定分布式内存模型，各进程通过相互发送消息交换信息。
应用使用这些通信 API 时，无须处理互连网络的具体细节。MPI 实现允许进程用逻辑编号彼此
寻址，很像电话系统使用电话号码：电话用户可以拨号联系对方，无须知道对方实际位于何处，
也无须知道通话经过怎样的路由。

\begin{figure}[htbp]
  \centering
  \begin{tikzpicture}[>=Stealth,node distance=0.35cm]
    \foreach \n in {0,...,3} {
      \begin{scope}[shift={(\n*3.4,0)}]
        \draw[rounded corners=3pt,thick,blue!60!black] (0,0) rectangle (2.8,2.15);
        \node[font=\small\bfseries] at (1.4,0.28) {节点 \n};
        \foreach \x/\y/\r in {0.65/1.55/0,1.55/1.55/1,0.65/0.85/2,1.55/0.85/3} {
          \node[draw,ellipse,fill=red!18,minimum width=0.82cm,minimum height=0.48cm,
                font=\scriptsize] at (\x,\y) {进程};
        }
      \end{scope}
    }
    \draw[thick] (1.4,-0.25) -- (11.6,-0.25);
    \foreach \x in {1.4,4.8,8.2,11.6} {\draw[thick] (\x,0) -- (\x,-0.25);}
    \node[below,font=\small] at (6.5,-0.25) {互连网络};
  \end{tikzpicture}
  \caption{程序员看到的 MPI 进程。依据原书图 23.4 重绘；每个节点可以包含一个或多个进程。}
  \label{fig:23-4}
\end{figure}

典型 MPI 应用会在进程之间划分数据和工作。如图 \ref{fig:23-4} 所示，每个节点可以包含
一个或多个进程，图中以节点内部的椭圆表示。进程向前执行时，可能需要其他进程的数据，例如
Jacobi 迭代法执行 halo 交换时就是如此；这种需求通过发送和接收消息满足。有时，进程还需要
彼此同步并产生集体结果，例如归约 $L_2$ 范数；这类操作由集体通信 API 函数完成。

MPI 程序与 CUDA 一样，都建立在 SPMD 并行编程模型之上：所有 MPI 进程执行同一个程序。
MPI 系统提供一组 API，用来建立允许进程相互通信的通信系统。图 \ref{fig:23-5} 列出为 MPI
应用建立和关闭通信系统所需的四类基本操作，图 \ref{fig:23-6} 则展示 MPI 应用通常如何
使用这些函数。

\noindent\textbf{译者注：}底本正文将图 \ref{fig:23-5} 概括为“five essential MPI
functions”，但原图只列出 \code{MPI_Init()}、\code{MPI_Comm_rank()}、
\code{MPI_Comm_size()} 和 \code{MPI_Finalize()} 四个函数；译稿依照原图和后续示例改为
“四类基本操作”。

\begin{figure}[htbp]
  \centering
  \begin{lstlisting}[language=C,numbers=none,basicstyle=\ttfamily\scriptsize]
int MPI_Init(int *argc, char ***argv)              // 初始化 MPI
int MPI_Comm_rank(MPI_Comm comm, int *rank)       // 返回 comm 中调用进程的秩
int MPI_Comm_size(MPI_Comm comm, int *size)       // 返回 comm 中的进程数
int MPI_Finalize()                                // 结束 MPI 应用并释放资源
  \end{lstlisting}
  \caption{建立和关闭通信系统所用的基本 MPI API 函数。依据原书图 23.5 逐字符重排，
  说明文字译为中文。}
  \label{fig:23-5}
\end{figure}

MPI 应用由 \code{mpirun} 或 \code{mpiexec} 启动。用户把程序的可执行文件交给命令，
命令再创建执行这一文件的进程。\footnote{MPI 标准建议把 \code{mpiexec} 作为启动 MPI
进程的标准命令名；Open MPI 等实现同时提供含义相同的 \code{mpirun}。}每个进程首先调用
\code{MPI_Init()} 初始化 MPI
运行时（图 \ref{fig:23-6} 第 3 行）；该调用为运行应用的所有进程初始化通信系统。

\begin{figure}[htbp]
  \centering
  \begin{lstlisting}[language=C,basicstyle=\ttfamily\scriptsize,firstnumber=1]
#include "mpi.h"
int main(int argc, char *argv[]) {
    MPI_Init(&argc, &argv);
    int rank;
    MPI_Comm_rank(MPI_COMM_WORLD, &rank);
    int numRanks;
    MPI_Comm_size(MPI_COMM_WORLD, &numRanks);
    ...
    MPI_Finalize();
    return 0;
}
  \end{lstlisting}
  \caption{一个简单的 MPI 主程序。依据原书图 23.6 逐字符重排。}
  \label{fig:23-6}
\end{figure}

接下来调用 \code{MPI_Comm_rank()}（第 5 行），它向每个调用进程返回一个唯一编号，称为
进程的 \textbf{MPI 秩（MPI rank）}。进程得到的编号范围是 0 到进程数减 1。MPI 秩在通信
中唯一标识一个进程，相当于电话系统中的电话号码；对于 CUDA 线程，它类似表达式
\code{blockIdx.x*blockDim.x + threadIdx.x}。CUDA 线程最多可以组织成三个维度；同样，
MPI 秩也可以用 MPI 笛卡尔通信器组织成多个维度。

\code{MPI_Comm_rank()} 接收两个参数。第一个参数的 MPI 内置类型为 \code{MPI_Comm}，
它指定请求的范围，也就是由一个 \code{MPI_Comm} 变量标识的进程集合。\code{MPI_Comm}
类型的变量通常称为\textbf{通信器（communicator）}。\code{MPI_Comm} 和其他 MPI 内置
类型都定义在 \code{mpi.h} 头文件中（第 1 行）；所有使用 MPI 的程序文件都应包含该头文件。
MPI 应用可以创建一个或多个通信器，每个通信器都指定调用通信函数时应包含的 MPI 进程组。
调用 \code{MPI_Comm_rank()} 时，它在通信器参数指定的组内为每个进程分配唯一 ID。图
\ref{fig:23-6} 传入 \code{MPI_COMM_WORLD}；这是默认使用的通信器，表示包含运行应用的
全部 MPI 进程。\footnote{关于在应用中创建并使用多个通信器的细节，尤其是组内通信器
（intracommunicator）与组间通信器（intercommunicator）的定义和用法，请参阅 MPI
参考手册\cite{ch23-mpi}。}

\code{MPI_Comm_rank()} 的第二个参数是整数变量的指针，函数把返回的秩写入该变量。图
\ref{fig:23-6} 为此声明了变量 \code{rank}。函数返回后，\code{rank} 保存调用进程在给定
通信器所指定进程组中的唯一 ID。MPI 文献经常交替使用“MPI 进程”和“MPI 秩”两个说法。

随后调用 \code{MPI_Comm_size()}（第 7 行），返回通信器中运行的 MPI 进程总数。一个
通信器指定的 MPI 进程总数称为该通信器的\textbf{秩数（number of ranks）}。
\code{MPI_Comm_size()} 同样接收两个参数：第一个是指定请求范围的 \code{MPI_Comm}；图
\ref{fig:23-6} 传入 \code{MPI_COMM_WORLD}，因而范围是应用的全部进程。第二个参数是整数
变量的指针，函数把返回值写入该变量；图中为此声明了 \code{numRanks}。函数返回后，
\code{numRanks} 就是运行应用的 MPI 进程数。这个数值由用户执行 \code{mpirun} 或
\code{mpiexec} 时配置。通过 \code{MPI_Comm_size()} 查询 MPI 运行时，可以在确定应创建
多少个分区时避免对应用的启动方式作隐含假设。

然后，应用执行主要计算（第 8 行），本节后续将展开这部分代码。主要计算结束后，应用调用
\code{MPI_Finalize()} 通知 MPI 运行时，并释放分配给应用的全部 MPI 通信资源（第 9 行）。
最后，应用以 0 为返回值退出，表示没有发生错误（第 10 行）。

现在考察各进程执行的主要计算。每个进程知道自己的秩和应用的秩数，因此可以确定自己负责
建模系统的哪一个分区。接着，进程用熟悉的 \code{cudaMalloc} 为该分区分配 GPU 内存，用
\code{cudaMemset} 初始化内存，并启动一个简单内核，依据被建模系统的边界条件初始化分区的
边界元素。

下面省略这些初始化代码，重点考察反复执行模板计算、halo 交换和 $L_2$ 范数归约的主机端
主要计算代码。假设进程 $P$ 在主机 CPU 上执行图 \ref{fig:23-7} 的代码，并在连接该主机的
GPU 全局内存中分配输出数组和输入数组。两个数组都包含 $ny\times nx$ 个元素，如图
\ref{fig:23-3} 所示。数组的第一行，也就是元素 0 到 $nx-1$，由顶部邻居提供的 halo 元素
组成；类似地，最后一行，即元素 $(ny-1)nx$ 到 $ny\,nx-1$，由底部邻居提供的 halo 元素
组成。

\begin{figure}[htbp]
  \centering
  \begin{lstlisting}[language=C++,basicstyle=\ttfamily\scriptsize,firstnumber=1]
for(unsigned int iter = 0; l2norm > tol && iter < iter_max; ++iter) {

    // Reset L2 norm
    cudaMemset(l2norm_d, 0, sizeof(float));

    // Compute stencil
    launch_jacobi_kernel(output, input, l2norm_d, ny, nx);
    cudaDeviceSynchronize();

    // Halo exchange
    const int top = rank > 0 ? rank - 1 : (numRanks - 1);
    const int bottom = (rank + 1) % numRanks;
    MPI_Sendrecv(output + nx, nx, MPI_FLOAT, top, 0,
                 output + (ny - 1)*nx, nx, MPI_FLOAT, bottom, 0,
                 MPI_COMM_WORLD, MPI_STATUS_IGNORE);
    MPI_Sendrecv(output + (ny - 2)*nx, nx, MPI_FLOAT, bottom, 0,
                 output, nx, MPI_FLOAT, top, 0,
                 MPI_COMM_WORLD, MPI_STATUS_IGNORE);

    // Copy L2 norm to host
    cudaMemcpy(l2norm_h, l2norm_d, sizeof(float), cudaMemcpyDeviceToHost);

    // Reduce L2 norm
    MPI_Allreduce(l2norm_h, &l2norm, 1, MPI_FLOAT, MPI_SUM, MPI_COMM_WORLD);
    l2norm = std::sqrt(l2norm);

    std::swap(output, input);

}
  \end{lstlisting}
  \caption{基于 MPI 的多 GPU Jacobi 迭代法主机代码。依据原书图 23.7 逐字符重排。}
  \label{fig:23-7}
\end{figure}

图 \ref{fig:23-7} 给出基于 MPI 的多 GPU 主机代码主循环。只要 $L_2$ 范数仍高于最小
容差阈值，并且迭代次数没有超过允许的最大值，循环就会继续（第 1 行）。每次迭代首先调用
\code{cudaMemset} 把 $L_2$ 范数平方和初始化为 0（第 4 行）。随后调用一个函数（第 7 行），
启动图 \ref{fig:23-1} 实现的 Jacobi 内核。该内核执行一次模板计算：从输入数组读取网格点
旧值，把新值写入输出数组，同时更新 $L_2$ 范数平方和。每次迭代都以双缓冲方式交换这两个
数组（第 27 行）。代码等待 Jacobi 内核执行完毕（第 8 行），以保证开始 halo 交换前，新网格
点值已经计算完成。

下一步执行 halo 交换。MPI 进程需要与邻居通信，彼此发送新的边界网格点值。这样的通信可以
用 MPI \textbf{点对点通信（point-to-point communication）}完成，它在两个通信进程之间
交换消息。MPI 提供许多点对点通信例程；本例采用 \code{MPI_Sendrecv()}。理解该函数以前，
先考察 MPI 提供的更基础点对点通信例程。

最基本的点对点通信例程是 \code{MPI_Send()} 和 \code{MPI_Recv()}：源进程调用前者向目标
进程发送数据，目标进程调用后者从源进程接收数据。这类似电话系统中一方拨号、另一方接听。
由于通信双方都必须参与，通信才能发生，这种点对点通信称为\textbf{双边通信（two-sided
communication）}。

\begin{figure}[htbp]
  \centering
  \begin{lstlisting}[language=C,numbers=none,basicstyle=\ttfamily\scriptsize]
int MPI_Send(void *buf, int count, MPI_Datatype datatype,
             int dest, int tag, MPI_Comm comm)
// buf: 发送缓冲区起始地址（指针）
// count: 发送缓冲区中的元素数（非负整数）
// datatype: 每个发送元素的数据类型（MPI_Datatype）
// dest: 目标进程的秩（整数）
// tag: 消息标签（整数）
// comm: 通信器（句柄）
  \end{lstlisting}
  \caption{\code{MPI_Send()} 函数的签名。依据原书图 23.8 逐字符重排，参数说明译为中文。}
  \label{fig:23-8}
\end{figure}

图 \ref{fig:23-8} 给出 \code{MPI_Send()} 的签名。第一个参数指向要发送数据所在内存区域的
起始位置；第二个整数参数给出要发送的数据元素数；第三个参数是 MPI 内置类型
\code{MPI_Datatype}，表示 MPI 库所看到的各数据元素类型。\code{mpi.h} 定义了这一类型
可以保存的值，包括 \code{MPI_DOUBLE}（双精度浮点数）、\code{MPI_FLOAT}（单精度浮点数）、
\code{MPI_INT}（整数）和 \code{MPI_CHAR}（字符）。这些类型的确切大小取决于主机处理器上
相应 C 类型的大小；关于 MPI 类型更复杂的用法，请参阅 MPI 参考手册\cite{ch23-mpi}。
第四个参数给出目标进程的 MPI 秩；第五个参数提供标签，可以把同一进程发送的消息分类；
第六个参数是通信器，指定定义目标 MPI 秩的进程组。

\begin{figure}[htbp]
  \centering
  \begin{lstlisting}[language=C,numbers=none,basicstyle=\ttfamily\scriptsize]
int MPI_Recv(void *buf, int count, MPI_Datatype datatype,
             int source, int tag, MPI_Comm comm, MPI_Status *status)
// buf: 接收缓冲区起始地址（指针）
// count: 接收缓冲区允许接收的最大元素数（整数）
// datatype: 每个接收元素的数据类型（MPI_Datatype）
// source: 源进程的秩（整数）
// tag: 消息标签（整数）
// comm: 通信器（句柄）
// status: 状态对象（MPI_Status）
  \end{lstlisting}
  \caption{\code{MPI_Recv()} 函数的签名。依据原书图 23.9 逐字符重排，参数说明译为中文。}
  \label{fig:23-9}
\end{figure}

图 \ref{fig:23-9} 给出 \code{MPI_Recv()} 的签名。第一个参数指向接收数据应写入的内存区域；
第二个整数参数给出该函数允许接收的最大元素数；第三个参数指定各接收元素的
\code{MPI_Datatype}；第四个参数给出消息源的秩；第五个参数指定目标进程期望的标签值。
如果目标进程不想限定某个特定标签，可以使用 \code{MPI_ANY_TAG}，表示接收方愿意接受来自
该源进程、具有任意标签值的消息。

执行 halo 交换时，一个进程把自己的边界数据发送给一个邻居，同时还要从另一个邻居接收
halo 数据。例如，在图 \ref{fig:23-3} 中，当进程 1（即控制 GPU 1 的进程）把顶部边界
网格点值发送给进程 0 时，进程 2 也会把自己的顶部边界网格点值发送给进程 1。所有进程都要
同时把边界数据发送给顶部邻居，并从底部邻居接收 halo 数据。为此，MPI 提供
\code{MPI_Sendrecv()}，它同时向一个目标发送数据，并从一个源接收数据。该函数签名见图
\ref{fig:23-10}；它结合了 \code{MPI_Send()} 与 \code{MPI_Recv()}，因此不再逐一解释参数。

\begin{figure}[htbp]
  \centering
  \begin{lstlisting}[language=C,numbers=none,basicstyle=\ttfamily\scriptsize]
int MPI_Sendrecv(void *sendbuf, int sendcount, MPI_Datatype sendtype,
                 int dest, int sendtag, void *recvbuf, int recvcount,
                 MPI_Datatype recvtype, int source, int recvtag,
                 MPI_Comm comm, MPI_Status *status)
// sendbuf/sendcount/sendtype: 发送缓冲区、元素数和元素类型
// dest/sendtag: 目标进程的秩和发送标签
// recvbuf/recvcount/recvtype: 接收缓冲区、最大元素数和元素类型
// source/recvtag: 源进程的秩和接收标签
// comm: 通信器；status: 接收操作的状态对象
  \end{lstlisting}
  \caption{\code{MPI_Sendrecv()} 函数的签名。依据原书图 23.10 逐字符重排，参数说明译为中文。}
  \label{fig:23-10}
\end{figure}

也可以用 MPI 非阻塞通信达到同一目的：通过 \code{MPI_Isend} 发起非阻塞发送，通过
\code{MPI_Irecv} 发起非阻塞接收，再用 \code{MPI_Wait} 等待两项操作完成。这些非阻塞通信
函数能够表达更复杂的模式，例如同时向不同邻居发送数据并从它们接收数据。不过，对于本例
所需的并发通信模式，\code{MPI_Sendrecv()} 只需一次 API 调用，表达起来更方便。关于非阻塞
通信函数的详细用法，请参阅 MPI 文档\cite{ch23-mpi}。

回到图 \ref{fig:23-7}，特别考察执行 halo 交换的第 11--18 行。首先，各进程确定顶部邻居
（第 11 行）和底部邻居（第 12 行）。代码没有简单地计算 \code{rank - 1} 和
\code{rank + 1}，而是采用回绕策略：最顶端的秩把最底端的秩视为顶部邻居，反之亦然。
这种回绕策略常称为\textbf{周期边界条件（periodic boundary condition）}。它只模拟大型
系统的一小部分，却不会引入比例过高的边缘效应。周期边界条件需要对最顶端和最底端的秩作
特殊处理，因为它们本来分别没有可通信的顶部邻居和底部邻居。

接下来，两次 \code{MPI_Sendrecv} 调用完成 halo 交换。第一次调用（第 13--15 行）把从
\code{output + nx} 开始的顶部边界行发送给顶部邻居，并从底部邻居接收从
\code{output + (ny - 1)*nx} 开始的底部 halo 行。第二次调用（第 16--18 行）把从
\code{output + (ny - 2)*nx} 开始的底部边界行发送给底部邻居，并从顶部邻居接收从
\code{output} 开始的顶部 halo 行。所有情况下，通信元素数都是一行的大小 $nx$，数据类型
都是 \code{MPI_FLOAT}，它对应 C 和 C++ 的 \code{float} 类型。

需要特别注意，传给 \code{MPI_Sendrecv} 的输出和输入指针都指向 GPU 设备内存。CUDA 感知
MPI 允许把 GPU 设备内存指针传给 MPI 函数。如果 MPI 实现不具备 CUDA 感知能力，程序员就
必须先把数据从设备内存复制到主机内存，再对主机内存数组调用 MPI 通信函数。相比之下，
CUDA 感知 MPI 实现理解 CUDA 编程模型，允许通信函数直接作用于 CUDA 全局内存数组。让
发送和接收函数直接处理设备数组，可以消除源代码中的 \code{cudaMemcpy} 调用，从而减少
主机代码膨胀和额外通信延迟。大多数 MPI 实现，如 MPICH\cite{ch23-mpich}、Open
MPI\cite{ch23-openmpi} 和 MVAPICH2\cite{ch23-mvapich}，都支持 CUDA 感知并可在构建时启用。

halo 交换结束后，主机代码需要计算 $L_2$ 范数，判断系统是否仍有显著变化、是否需要继续
迭代。每个 GPU 已经计算了自己分区残差的 $L_2$ 范数平方和；现在还要跨所有分区把这些值
相加。换句话说，需要跨全部 MPI 秩执行归约，每个秩贡献自己分区的 $L_2$ 范数平方和。
这样的跨秩归约可以用 MPI 集体通信完成。

MPI 点对点通信只涉及两个秩之间的通信；与之不同，\textbf{MPI 集体通信（collective
communication）}涉及通信器中的全部秩。常用集体通信类型包括屏障、广播、归约、聚集和
散布。本章不逐一展开这些函数；细节请参阅 MPI 编程指南\cite{ch23-mpi}。一般来说，MPI
运行时开发者和系统供应商会高度优化集体通信函数。与自行组合多次发送和接收调用实现同一
功能相比，使用集体通信函数通常性能更好，代码也更易读、开发效率更高。

计算 $L_2$ 范数时，我们需要的 MPI 集体通信函数是 \code{MPI_Allreduce}。该函数从每个秩
取得一份数值缓冲区，把这些缓冲区归约成一个缓冲区，再把结果缓冲区的副本交给每个秩。图
\ref{fig:23-11} 给出函数签名。第一个参数 \code{sendbuf} 指定各秩贡献给归约操作的输入
缓冲区；第二个参数 \code{recvbuf} 指定最终结果缓冲区的存放位置。也就是说，接收缓冲区的
每个元素，都是所有秩发送缓冲区同一下标元素的归约结果。第三个参数 \code{count} 指定发送
缓冲区的元素数，也等于接收缓冲区的元素数；第四个参数 \code{datatype} 指定整数、浮点数
等被归约元素的类型；第五个参数 \code{op} 指定归约操作。MPI 内置许多归约操作，包括求和、
求积、最小值、最大值和逻辑运算等。最后一个参数 \code{comm} 指定参与归约的通信器。

\begin{figure}[htbp]
  \centering
  \begin{lstlisting}[language=C,numbers=none,basicstyle=\ttfamily\scriptsize]
int MPI_Allreduce(const void *sendbuf, void *recvbuf, int count,
                  MPI_Datatype datatype, MPI_Op op, MPI_Comm comm)
// sendbuf: 发送缓冲区起始地址
// recvbuf: 接收缓冲区起始地址
// count: 发送缓冲区中的元素数
// datatype: 发送缓冲区元素的数据类型
// op: 归约操作；comm: 通信器
  \end{lstlisting}
  \caption{\code{MPI_Allreduce()} 函数的签名。依据原书图 23.11 逐字符重排，参数说明译为中文。}
  \label{fig:23-11}
\end{figure}

再次回到图 \ref{fig:23-7}，考察归约 $L_2$ 范数的第 20--25 行。代码先把本分区的 $L_2$
范数平方和从设备复制到主机（第 21 行），再用 \code{MPI_Allreduce} 跨所有秩归约这些平方和
（第 24 行）。这里缓冲区中的元素数是 1，因为每个秩只贡献一个数值，即本分区的 $L_2$
范数平方和。归约操作是 \code{MPI_SUM}，表示求所有秩贡献值之和。\code{MPI_Allreduce}
返回后，每个秩都得到全部 $L_2$ 范数平方和的总和；代码对它开平方，得到整个系统最终的
$L_2$ 范数（第 25 行），并用它判断循环是否应继续（第 1 行）。

读者可能会问，为什么要在调用 \code{MPI_Allreduce} 以前，把 $L_2$ 范数平方和从设备复制
到主机。既然 MPI 具备 CUDA 感知能力，原本可以直接对 GPU 缓冲区调用
\code{MPI_Allreduce}。原因在于，主机需要归约结果来检查是否收敛，所以无论如何都必须把
结果复制到主机。在这种情况下，先复制到主机，再调用 \code{MPI_Allreduce} 更高效。

\section{重叠计算与通信}
\label{sec:overlap-computation-communication}

第 \ref{sec:multi-gpu-mpi} 节的多 GPU 模板实现中，每个进程先计算整个分区的新值，再与顶部
和底部邻居交换 halo 数据，然后重复这一过程。这种策略十分简单，却不够高效，因为它迫使
系统在两种模式之间切换，而每种模式只利用系统的一部分资源。第一种模式下，所有进程都用
GPU 执行计算，通信网络闲置；第二种模式下，所有进程都通过通信网络与顶部、底部邻居交换
halo 数据，GPU 计算硬件利用不足。理想情况下，应当让通信网络与 GPU 计算硬件始终同时工作。
如图 \ref{fig:23-12} 所示，把每个进程的计算任务分成两个阶段，就可以提高利用率。

\begin{figure}[htbp]
  \centering
  \begin{tikzpicture}[x=0.50cm,y=0.46cm,>=Stealth]
    \node[font=\small\bfseries] at (3.5,9.0) {(a) 阶段 1：计算边界值};
    \node[font=\small\bfseries,align=center] at (13,9.0) {(b) 阶段 2：计算内部值，\\同时执行 halo 交换};
    \foreach \base in {0,9.5} {
      \foreach \y/\c/\g in {5.8/red!68/0,3.3/yellow!75/1,0.8/green!60/2} {
        \draw[thick] (\base,\y) rectangle +(7,1.7);
        \draw[dashed] (\base,\y+0.17) -- +(7,0);
        \draw[dashed] (\base,\y+1.53) -- +(7,0);
        \node[left,font=\scriptsize] at (\base,\y+0.85) {GPU \g};
      }
    }
    \fill[red!68] (0,5.97) rectangle +(7,0.2);
    \fill[yellow!75] (0,3.47) rectangle +(7,0.2);
    \fill[yellow!75] (0,4.83) rectangle +(7,0.2);
    \fill[green!60] (0,0.97) rectangle +(7,0.2);
    \fill[green!60] (0,2.33) rectangle +(7,0.2);
    \foreach \y/\c in {5.8/red!68,3.3/yellow!75,0.8/green!60}
      \fill[\c] (9.5,\y) rectangle +(7,1.7);
    \foreach \y in {5.8,3.3} {
      \draw[<->,orange!80!black] (16.7,\y-0.03) .. controls (17.3,\y-0.35) and
        (17.3,\y-0.65) .. (16.7,\y-0.87);
      \node[right,font=\scriptsize,align=left] at (17.05,\y-0.45) {halo\\交换};
    }
  \end{tikzpicture}
  \caption{让计算与通信重叠的两阶段策略。依据原书图 23.12 重绘。}
  \label{fig:23-12}
\end{figure}

第一阶段中，每个进程计算下一次迭代时相邻进程需要作为 halo 点的边界行。图
\ref{fig:23-12} 用着色行表示本阶段计算的边界行。例如，GPU 1 计算自己的顶部和底部边界
行，它们分别是 GPU 0 与 GPU 2 的 halo 行。之所以先计算这些边界行，是因为相邻进程在下次
迭代中需要它们；这样，第二阶段就可以一边把边界行传给邻居，一边计算其余内部网格点。

第二阶段中，每个进程并行执行两项活动。第一项是把新边界值传给相邻进程；第二项是计算分区
中的其余数据。如果通信时间短于计算时间，就可以隐藏通信延迟，并在绝大多数时间内充分利用
计算硬件。通常，只要每个分区的内部部分有足够多的行，使各进程能执行足够多的计算来遮蔽
通信，就可以做到这一点。

\begin{figure}[htbp]
  \centering
  \begin{tikzpicture}[x=0.72cm,y=0.72cm,>=Stealth,font=\scriptsize]
    \draw[->] (0,6.4) -- (12.4,6.4) node[above left] {时间};
    \node[left,align=right] at (0,5.75) {计算硬件\\（GPU）};
    \node[left,align=right] at (0,4.65) {通信网络};
    \draw[fill=blue!30] (0,5.35) rectangle (6.9,6.05) node[midway] {计算全部网格点};
    \draw[fill=orange!28] (7.0,4.25) rectangle (9.25,5.0) node[midway,align=center] {顶部 halo\\交换};
    \draw[fill=green!25] (9.25,4.25) rectangle (11.5,5.0) node[midway,align=center] {底部 halo\\交换};
    \draw[fill=blue!25] (11.5,4.25) rectangle (12.35,5.0) node[midway,align=center] {归约\\$L_2$};
    \node[align=center] at (6.2,3.75) {(a) 计算与通信不重叠};

    \draw[->] (0,3.1) -- (12.4,3.1) node[above left] {时间};
    \node[left,align=right] at (0,2.25) {计算硬件\\（GPU）};
    \node[left,align=right] at (0,0.85) {通信网络};
    \draw[fill=blue!30] (0,2.35) rectangle (7.0,2.95) node[midway] {计算内部网格点};
    \draw[fill=orange!35] (0,1.75) rectangle (1.75,2.35) node[midway,font=\tiny] {顶部边界};
    \draw[fill=green!35] (0,1.15) rectangle (1.75,1.75) node[midway,font=\tiny] {底部边界};
    \draw[fill=orange!28] (1.75,0.45) rectangle (4.5,1.1) node[midway,align=center] {顶部 halo 交换};
    \draw[fill=green!25] (4.5,0.45) rectangle (7.0,1.1) node[midway,align=center] {底部 halo 交换};
    \draw[fill=blue!25] (7.0,0.45) rectangle (8.0,1.1) node[midway,align=center] {归约\\$L_2$};
    \draw[->,orange!80!black] (1.75,2.05) -- (1.75,1.1);
    \draw[->,green!60!black] (1.75,1.45) -- (4.5,1.1);
    \node[align=center] at (6.2,-0.05) {(b) 计算与通信重叠};
  \end{tikzpicture}
  \caption{计算与通信不重叠和重叠时的执行时间线比较。依据原书图 23.13 重绘。}
  \label{fig:23-13}
\end{figure}

图 \ref{fig:23-13} 比较计算与通信不重叠和重叠时的执行时间线。在图
\ref{fig:23-13}(a) 中，计算不与通信重叠：计算硬件执行全部计算时通信网络闲置；随后通信
网络完成顶部和底部 halo 交换时，计算硬件闲置；最后还要执行 $L_2$ 范数归约。图
\ref{fig:23-13}(b) 则把顶部、底部网格点与内部网格点分开计算，相应边界点一旦算出，就立即
发起顶部或底部 halo 交换。GPU 计算资源足够时，三类网格点可以并发计算，但应优先计算边界
网格点，以便尽早发起 halo 交换。关键在于，halo 交换期间仍继续计算内部网格点；计算硬件与
通信网络因而同时得到利用，总执行时间缩短。不可避免的是，归约 $L_2$ 范数依赖所有网格点
都已算完，因此无法与这些计算重叠。

实现这种重叠，需要把一次 Jacobi 内核调用拆成分别处理边界元素和内部元素的多个内核调用。
例如，可以拆成三个调用：一个计算顶部边界行，一个计算底部边界行，另一个计算内部元素。
相应边界行内核一结束，就可以发起 halo 交换。为了让三个内核并发执行，并分别让通信活动与
各内核独立同步，需要使用 CUDA 编程模型的重要特性：\textbf{CUDA stream}。

CUDA stream 是按顺序排列的 CUDA 操作序列，操作可以包括内核调用、内存复制等。同一
stream 中的全部操作严格按照放入顺序依次执行；不同 stream 中的操作则没有彼此之间的顺序
约束，可以并发执行。因此，要让上述三个 Jacobi 内核并发执行，应把它们放入三个不同的
stream。首先创建 stream 对象：

\begin{lstlisting}[language=C++,numbers=none]
cudaStream_t topStream, bottomStream, internalStream;
cudaStreamCreate(&internalStream);
cudaStreamCreate(&topStream);
cudaStreamCreate(&bottomStream);
\end{lstlisting}

第一行声明三个 \code{cudaStream_t} 类型的 stream 对象，分别命名为 \code{topStream}、
\code{bottomStream} 和 \code{internalStream}，对应处理顶部边界行、底部边界行和内部行的
内核调用。接下来三行分别调用 \code{cudaStreamCreate}，传入 stream 对象的指针；每次调用
都会创建一个新 stream，并初始化它的状态。

要把内核调用放入某个 stream，可以把 stream 对象作为内核调用的执行配置参数。第 2 章已经
介绍，内核调用的尖括号中有两个必需的执行配置参数：线程网格维度和线程块维度。第 5 章还
介绍了可选的第三个参数，它指定动态分配的共享内存大小，省略时默认为 0。类似地，内核调用
还接受可选的第四个参数，指定应把调用放入哪个 stream；省略时调用会进入默认 stream。例如：

\begin{lstlisting}[language=C++,numbers=none]
kernelName<<<dimGrid, dimBlock, 0, myStream>>>(...)
\end{lstlisting}

这段代码调用内核 \code{kernelName}，线程网格维度由 \code{dimGrid} 指定，线程块维度由
\code{dimBlock} 指定，动态分配的共享内存为 0，并把调用放入 \code{myStream}。注意，
调用返回并不表示内核已经执行，只表示内核已经放入指定 stream。内核一进入 stream，主机
代码就会继续执行后面的语句。CUDA 内核调用之所以是异步的，正是因为主机端调用只是把内核
放入 stream，而不是亲自执行内核。主机代码继续执行也不表示相应线程网格已经完成内核，
甚至不表示线程网格已经开始执行；它只表示内核已经进入指定 stream，等该 stream 中所有
先前操作执行完毕后，最终会启动线程网格执行内核。

\begin{figure}[htbp]
  \centering
  \begin{lstlisting}[language=C++,basicstyle=\ttfamily\scriptsize,firstnumber=1]
for(unsigned int iter = 0; l2norm > tol && iter < iter_max; ++iter) {

    // Reset L2 norm
    cudaMemsetAsync(l2norm_d, 0, sizeof(float), internalStream);
    cudaEventRecord(resetL2, internalStream);
    cudaStreamWaitEvent(topStream, resetL2, 0);
    cudaStreamWaitEvent(bottomStream, resetL2, 0);

    // Compute stencil
    launch_jacobi_kernel(output, input, l2norm_d, 3, nx, topStream);
    cudaEventRecord(computeTop, topStream);
    launch_jacobi_kernel(output + (ny - 3)*nx, input + (ny - 3)*nx, l2norm_d, 3,
                         nx, bottomStream);
    cudaEventRecord(computeBottom, bottomStream);
    launch_jacobi_kernel(output + nx, input + nx, l2norm_d, ny - 2, nx,
                         internalStream);

    // Copy L2 norm to host
    cudaStreamWaitEvent(internalStream, computeTop, 0);
    cudaStreamWaitEvent(internalStream, computeBottom, 0);
    cudaMemcpyAsync(l2norm_h, l2norm_d, sizeof(float), cudaMemcpyDeviceToHost,
                    internalStream);

    // Halo exchange
    const int top = rank > 0 ? rank - 1 : (numRanks - 1);
    const int bottom = (rank + 1) % numRanks;
    cudaStreamSynchronize(topStream);
    MPI_Sendrecv(output + nx, nx, MPI_FLOAT, top, 0,
                 output + (ny - 1)*nx, nx, MPI_FLOAT, bottom, 0,
                 MPI_COMM_WORLD, MPI_STATUS_IGNORE);
    cudaStreamSynchronize(bottomStream);
    MPI_Sendrecv(output + (ny - 2)*nx, nx, MPI_FLOAT, bottom, 0,
                 output, nx, MPI_FLOAT, top, 0,
                 MPI_COMM_WORLD, MPI_STATUS_IGNORE);

    // Reduce L2 norm
    cudaStreamSynchronize(internalStream);
    MPI_Allreduce(l2norm_h, &l2norm, 1, MPI_FLOAT, MPI_SUM, MPI_COMM_WORLD);
    l2norm = std::sqrt(l2norm);

    std::swap(output, input);

}
  \end{lstlisting}
  \caption{基于 MPI、重叠计算与通信的多 GPU Jacobi 迭代法主机代码。依据原书图 23.14
  逐字符重排。}
  \label{fig:23-14}
\end{figure}

了解如何创建 stream 并把内核调用放入其中以后，现在把 Jacobi 内核拆成三个调用，分别放入
三个 stream。图 \ref{fig:23-14} 给出基于 MPI、重叠计算与通信的多 GPU Jacobi 迭代法
主机代码。图 \ref{fig:23-7} 每次迭代只调用一次 \code{launch_jacobi_kernel}、启动一个内核；
相比之下，图 \ref{fig:23-14} 分别调用它计算顶部边界行（第 10 行）、底部边界行（第 12--13
行）和内部网格点（第 15--16 行）。

计算顶部边界行时（第 10 行），函数需要前三行作为输入，也就是顶部 halo 行、顶部边界行和
顶部内部行。因此，代码原样传入 \code{output} 与 \code{input} 指针，并以 3 作为第 4 个参数，
表示关心三行输入。重要的是，还要传入 \code{topStream}，\code{launch_jacobi_kernel} 会把
内核调用放入这一 stream。

计算底部边界行时（第 12--13 行），函数需要最后三行作为输入，也就是底部 halo 行、底部边界
行和底部内部行。因此，代码传入指向最后三行开头的 \code{output + (ny - 3)*nx} 与
\code{input + (ny - 3)*nx}，并以 3 作为第 4 个参数；还传入 \code{bottomStream}，作为
内核调用所在的 stream。

最后，计算内部网格点时（第 15--16 行），函数需要除顶部和底部 halo 行以外的所有行，因此
传入跳过第一行的 \code{output + nx} 与 \code{input + nx}，并传入 \code{ny - 2}，排除两条
halo 行；还传入 \code{internalStream}，作为内核调用所在的 stream。

在三个 stream 中并发执行三个内核时，有一点非常重要：这些内核可能以任意顺序执行。最可能
的情况是，计算边界网格点的内核因为先被调用而先执行，但不能保证这一顺序。计算内部网格点
的内核可能抢先执行并占满 GPU 资源，使计算边界点的内核及随后的 halo 交换无法执行，从而
达不到重叠计算与通信的目的。为避免这种情况，需要让计算边界点的内核具有比计算内部点的
内核更高的优先级；也就是让 \code{topStream} 和 \code{bottomStream} 的优先级高于
\code{internalStream}。

默认情况下，\code{cudaStreamCreate} 创建的 stream 优先级相同。要为不同 stream 指定不同
优先级，需要使用其他 CUDA API：

\begin{lstlisting}[language=C++,numbers=none]
int low, high;
cudaDeviceGetStreamPriorityRange(&low, &high);
cudaStreamCreateWithPriority(&internalStream, cudaStreamDefault, low);
cudaStreamCreateWithPriority(&topStream, cudaStreamDefault, high);
cudaStreamCreateWithPriority(&bottomStream, cudaStreamDefault, high);
\end{lstlisting}

第二行查询 CUDA 运行时支持的 stream 优先级范围，也就是最高值和最低值；所用函数为
\texttt{cudaDeviceGetStream\allowbreak PriorityRange}。随后三行调用
\texttt{cudaStreamCreate\allowbreak WithPriority}。代码以最低优先级创建
\code{internalStream}，再以最高
优先级分别创建 \code{topStream} 和 \code{bottomStream}。本例只需区分两类 stream，因此
两个优先级值就足够；其他场景可能需要
两个以上的优先级层次，也可以使用中间值。优先级的绝对数值并不重要，只有相对关系有意义。

把三个 Jacobi 内核调用放入不同 stream 以后，它们可以并发执行。下面考察如何等待某个特定
内核结束，再发起 halo 交换等后续任务。为此，CUDA 提供
\code{cudaStreamSynchronize}。该函数会等待某个特定 stream 中此前已有的全部任务执行完毕。
如图 \ref{fig:23-14} 所示，代码先等待顶部边界行内核在 \code{topStream} 中结束（第 27 行），
再调用 \code{MPI_Sendrecv} 交换顶部边界行。类似地，代码等待底部边界行内核在
\code{bottomStream} 中结束（第 31 行），再调用 \code{MPI_Sendrecv} 交换底部边界行。
在此期间，由于没有对包含内部行内核的
\code{internalStream} 执行同步，该内核可以在两次 MPI halo 交换期间继续执行。这样就让
内部行计算与 halo 行通信重叠起来。

\noindent\textbf{译者注：}底本在解释底部边界同步时写作“第 32 行”，但图
\ref{fig:23-14} 第 31 行才是对 \code{bottomStream} 的同步调用；第 32 行已经是
\code{MPI_Sendrecv}。译稿按代码实际行号作最小修正。

虽然把一次内核调用拆成三个调用，可以为 halo 交换分别同步各个调用，但这也给 $L_2$ 范数
计算带来挑战。三个内核都会更新同一个 $L_2$ 范数平方和，所以它们都必须等这个值初始化为
0 后才能执行，并且复制到主机、跨秩归约之前，三个内核又必须全部完成。然而，三个内核调用
位于三个不同的 stream，而初始化与复制 $L_2$ 范数平方和的调用只能放入一个 stream。要解决
这一问题，必须让一个 stream 中的操作等待另一个 stream 中的操作完成；这种 stream 间交互
可以用 CUDA 事件实现。

CUDA \textbf{事件（event）}是标记 stream 中某个位置的标记。CUDA 允许程序员对特定事件
执行同步，也就是等待 stream 的执行到达事件所标记的位置。主机线程可以执行这种同步，也
可以把对某个事件的等待操作放入另一个 stream。初始化和复制 $L_2$ 范数时，需要的是后者。

图 \ref{fig:23-7} 使用阻塞函数 \code{cudaMemset} 把 GPU 全局内存中的 $L_2$ 范数平方和
初始化为 0；只有初始化完成，该函数才返回。图 \ref{fig:23-14} 改用非阻塞变体
\code{cudaMemsetAsync}（第 4 行），它与内核调用一样，接收额外参数并把初始化操作放入指定
stream。本例把它放入 \code{internalStream}，从而保证该 stream 中后续的内部行内核调用不会
在初始化结束前执行。

不过，另外两个计算顶部、底部边界行的内核也必须等 \code{cudaMemsetAsync} 完成。为此，
可以在 \code{internalStream} 的 \code{cudaMemsetAsync} 之后放入事件，再在
\code{topStream} 和 \code{bottomStream} 中、两个内核之前放入等待这一事件的同步操作。
下面的代码创建用于这种同步的事件：

\begin{lstlisting}[language=C++,numbers=none]
cudaEvent_t resetL2;
cudaEventCreateWithFlags(&resetL2, cudaEventDisableTiming);
\end{lstlisting}

第一行声明 \code{cudaEvent_t} 类型的事件对象，第二行创建该事件。第一个参数是事件对象的
指针；第二个参数
\texttt{cudaEventDisable\allowbreak Timing} 是编译期常量，表示不把
该事件用于计时；运行时因而可以更高效地实现基于事件的同步。事件创建以后，图
\ref{fig:23-14} 第 5 行调用 \code{cudaEventRecord}，把它放在 \code{internalStream} 的
\code{cudaMemsetAsync} 之后。此时，事件 \code{resetL2} 标记初始化操作的完成位置。

要在另一条 stream 中插入对该事件的等待操作，可以像图 \ref{fig:23-14} 第 6--7 行那样调用
\code{cudaStreamWaitEvent}。第一个参数指定插入等待操作的 stream，即第 6 行的
\code{topStream} 和第 7 行的 \code{bottomStream}；第二个参数指定等待的事件，二者都是
\code{resetL2}；第三个参数决定是否使用特殊选项，取 0 表示采用默认选项。

接着考察三个内核都结束以后，如何复制 $L_2$ 范数平方和。图 \ref{fig:23-7} 使用阻塞的
\code{cudaMemcpy}；图 \ref{fig:23-14} 改用非阻塞变体 \code{cudaMemcpyAsync}（第 21 行），
后者还接受指定目标 stream 的额外参数。本例把它放入 \code{internalStream}，因此它会等待
该 stream 中先前的内部行内核。

复制还必须等待另外两个边界行内核。为此，可以在 \code{topStream} 和
\code{bottomStream} 的内核调用后记录事件，再在 \code{internalStream} 的
\code{cudaMemcpyAsync} 之前插入对这些事件的等待操作。图 \ref{fig:23-14} 正是如此：第 11
行和第 14 行调用 \code{cudaEventRecord}，分别在 \code{topStream} 与 \code{bottomStream}
中记录事件；第 19--20 行在内存复制操作前，把等待两个事件的操作放入
\code{internalStream}。

图 \ref{fig:23-14} 还利用 \code{cudaMemcpyAsync} 非阻塞这一事实作了进一步优化：如图中
所示，在调用 \code{MPI_Sendrecv} 以前，就把从全局内存向主机内存复制 \code{l2norm} 的
操作放入 stream。如果计算内部网格点的内核比 halo 交换先结束，这会使 GPU 到主机 CPU 的
内存复制与网络 halo 交换重叠。不过，在调用 \code{MPI_Allreduce} 以前，必须保证内存复制
已经完成；第 37 行调用 \code{cudaStreamSynchronize}，等待执行 \code{cudaMemcpyAsync}
的 \code{internalStream}，从而提供这一保证。

\begin{figure}[htbp]
  \centering
  \resizebox{\linewidth}{!}{%
  \begin{tikzpicture}[x=1cm,y=1cm,>=Stealth,font=\scriptsize]
    \draw[->] (0,6.4) -- (16.8,6.4) node[above left] {时间};
    \foreach \y/\name in {5.5/{主机},4.0/{顶部 stream},2.5/{底部 stream},1.0/{内部 stream}}
      \node[left,align=right] at (0,\y) {\name};
    \foreach \x/\n in {0/04,0.7/05,1.4/06,2.1/07,2.8/10,3.5/11,4.2/12,4.9/14,5.6/15,6.3/19,7/20,7.7/21}
      \node[draw,minimum width=0.7cm,minimum height=0.5cm] at (\x+0.35,5.5) {\n};
    \node[draw,minimum width=3.2cm,minimum height=0.65cm] at (10.0,5.5) {MPI 交换顶部};
    \node[draw,minimum width=3.2cm,minimum height=0.65cm] at (13.3,5.5) {MPI 交换底部};
    \node[draw,minimum width=2.2cm,minimum height=0.65cm] at (16.1,5.5) {MPI 全归约};
    \node[draw,dotted,minimum width=1.7cm] at (2.1,4.0) {等待 resetL2};
    \node[draw,minimum width=2.0cm] at (4.0,4.0) {Jacobi（顶部）};
    \node[draw,dotted,minimum width=1.2cm] at (6.0,4.0) {computeTop};
    \node[draw,dotted,minimum width=1.7cm] at (2.1,2.5) {等待 resetL2};
    \node[draw,minimum width=2.0cm] at (4.0,2.5) {Jacobi（底部）};
    \node[draw,dotted,minimum width=1.5cm] at (6.2,2.5) {computeBottom};
    \node[draw,minimum width=1.7cm] at (1.0,1.0) {memset};
    \node[draw,dotted,minimum width=1.2cm] at (2.55,1.0) {resetL2};
    \node[draw,minimum width=6.0cm] at (7.1,1.0) {Jacobi（内部）};
    \node[draw,dotted,minimum width=0.8cm] at (10.75,1.0) {等待顶};
    \node[draw,dotted,minimum width=0.8cm] at (12.0,1.0) {等待底};
    \node[draw,minimum width=1.4cm] at (13.4,1.0) {memcpy};
    \draw[dotted,->] (2.55,1.35) -- (2.1,3.65);
    \draw[dotted,->] (2.55,1.35) -- (2.1,2.85);
    \draw[dotted,->] (6.0,3.65) -- (10.75,1.35);
    \draw[dotted,->] (6.2,2.15) -- (12.0,1.35);
    \draw[dashed,thick,->] (6.0,4.35) -- (8.4,5.15);
    \draw[dashed,thick,->] (6.2,2.85) -- (11.7,5.15);
    \draw[dashed,thick,->] (13.4,1.35) -- (15.2,5.15);
  \end{tikzpicture}
  }
  \caption{使用 CUDA stream 和事件重叠计算与 MPI 通信的示例时间线。依据原书图 23.15
  重绘；实线框为计算或复制操作，点线框为事件与等待，粗虚线连接 stream 进度与主机同步。}
  \label{fig:23-15}
\end{figure}

图 \ref{fig:23-15} 说明图 \ref{fig:23-14} 的主机代码如何在 Jacobi 求解器每次迭代的第
4--41 行中，重叠内部网格点计算与 halo 交换。四条时间线分别表示主机 CPU、
\code{topStream}、\code{bottomStream} 和 \code{internalStream} 的活动。主机 CPU 时间线
中的每个方格，要么是向某条 stream 插入操作（用图 \ref{fig:23-14} 的主机代码行号标记），
要么表示主机 CPU 自己执行的活动。

对于某条 stream 中的每个操作，主机时间线上的方格表示插入时刻，stream 时间线上的方格表示
GPU 执行时段；实线把插入动作与执行动作关联起来。例如，主机时间线上第 10 行的方格把计算
顶部网格点的 Jacobi 内核放入 \code{topStream}。这个内核必须等待 \code{resetL2} 事件，
所以执行会被推迟；如果复位 \code{l2norm} 所需时间更短或更长，Jacobi 内核就会相应地更早
或更晚执行。

点线箭头表示事件如何允许某个等待操作完成。例如，\code{resetL2} 事件有一条点线箭头指向
\code{topStream} 中对它的等待操作。\code{cudaMemset} 一结束，stream 就到达
\code{resetL2} 事件，等待操作随之完成，并允许同一 stream 中计算顶部网格点的 Jacobi 内核
继续执行。

粗虚线表示某条 stream 到达一个位置以后，主机上相应的同步调用便可以返回，主机因而能够
继续后续活动；这里的同步函数是 \code{cudaStreamSynchronize}。例如，图中
\code{bottomStream} 的
\code{computeBottom} 事件之后有一条粗虚线，表示主机何时可以从针对 \code{bottomStream}
的同步调用返回。由于该同步函数在此位置之后才开始执行，它会立即返回，解除对底部网格点
halo 交换的阻塞。

\noindent\textbf{译者注：}底本本段把 API 名写作 \code{cudaStreamsynchronize}；图
\ref{fig:23-14} 的代码及 CUDA API 正确名称均为 \code{cudaStreamSynchronize}。译稿按正确
符号作最小修正。

图 \ref{fig:23-15} 还揭示了一个重要现象：主机同步可能在计算操作完成（例如顶部 stream 的
Jacobi 内核）与依赖它的数据传输操作开始（例如主机执行的顶部 MPI 交换）之间引入显著延迟。
这项额外延迟通常称为\textbf{同步开销（synchronization overhead）}，它会显著影响操作之间
实际获得的重叠程度，例如内部 stream 的 Jacobi 内核与主机 MPI 交换之间的重叠。由于 MPI
交换调用是阻塞的，主机只有等顶部 MPI 交换完成以后，才能开始同步底部 stream 并执行底部
MPI 交换。这会在底部 stream 的 Jacobi 内核结束与底部单元的 MPI 交换开始之间增加延迟，
减少 MPI 交换与内部 stream Jacobi 内核的重叠。降低这种同步开销，是第
\ref{sec:multi-gpu-nccl} 节 NCCL 和第 \ref{sec:multi-gpu-nvshmem} 节 NVSHMEM 的重要设计目标。

还有一个重要细节：\code{l2norm_h} 指向的主机内存参与 \code{cudaMemcpyAsync} 等异步内存
操作，因此不能用 \code{malloc()} 等普通主机内存分配方法取得，而必须分配为
\textbf{固定页内存（pinned memory）}，有时也称为页锁定内存（page-locked memory）。
可以用 \code{cudaMallocHost} 分配固定页内存：

\begin{lstlisting}[language=C++,numbers=none]
float* l2norm_h;
cudaMallocHost(&l2norm_h, sizeof(float));
\end{lstlisting}

要完整理解固定页内存及异步复制为何需要它，先补充一些操作系统内存管理背景。现代计算机
系统中，操作系统同时管理许多应用的内存地址空间。每个应用都能访问一个庞大、逻辑上连续的
\textbf{虚拟地址空间（virtual address space）}。实际上，系统的物理内存容量有限，必须由
所有运行中的应用共享。为实现这种共享，操作系统把应用的虚拟地址空间划分成虚拟页，再把
这些虚拟页映射到物理内存中的物理页。

虚拟页到物理页的映射实现了多个目标。其一是保证进程若未作显式共享安排，就不能访问其他
进程在物理内存中的页面；进程只能访问映射到自己虚拟地址空间的物理页。其二是在内存需求
超过可用物理内存时仍能运行应用：只有最活跃的虚拟页映射到物理页，其余虚拟页放在磁盘上。
在这一过程中，操作系统可能需要把物理内存中的某些页面“换出（page out）”到磁盘等大容量
存储，为其他页面腾出空间。因此，应用的数据在执行期间随时可能被换出物理内存。

\code{cudaMemcpy()} 的实现会使用一种称为\textbf{直接内存访问（Direct Memory Access，
DMA）}设备的硬件。在主机与设备内存之间复制时，\code{cudaMemcpy()} 利用 DMA 完成任务。
在主机内存一侧，DMA 硬件使用物理地址，所以操作系统必须向 DMA 提供转换后的物理地址。
然而，DMA 操作结束前，数据有可能被换出；相应物理内存位置可能重新分配给另一虚拟内存位置
的数据。这样，分页活动就可能覆盖 DMA 正在处理的数据，使复制损坏。

常见的解决办法是让 CUDA 运行时分两步执行复制。对于主机到设备复制，运行时先把源主机
内存数据复制到“固定”的内存缓冲区，也就是把这些内存位置标记为分页机制不得换出；然后由
DMA 设备把固定页缓冲区复制到设备内存。对于设备到主机复制，运行时先由 DMA 设备把数据
从设备内存复制到固定页缓冲区，再把缓冲区复制到目标主机内存位置。借助额外的固定页内存
缓冲区，DMA 复制不会受到分页活动干扰。

这种方法有两个问题。第一，额外一次复制会增加 \code{cudaMemcpy()} 的延迟。第二，增加的
复杂性会使 \code{cudaMemcpy()} 以同步方式实现：只有函数完成并返回以后，主机才能继续执行，
因为主机必须参与内存复制。为了支持 \code{cudaMemcpyAsync()} 这样的异步复制，需要消除
主机把数据从普通缓冲区复制到固定页缓冲区的步骤。因此，\code{cudaMemcpyAsync()} 要求传入
的主机内存缓冲区本身就已经分配为固定页内存。

\section{使用 NCCL 的多 GPU 模板计算}
\label{sec:multi-gpu-nccl}

第 \ref{sec:overlap-computation-communication} 节的多 GPU Jacobi 迭代法通过并发的异步内核
调用、内存复制和 CUDA stream，有效重叠了计算与通信。不过，实现中仍有一些同步点会造成
低效。回看图 \ref{fig:23-14}，可以看到主机线程在调用 \code{MPI_Sendrecv} 前先调用
\code{cudaStreamSynchronize}（第 27、31 行），以保证计算顶部、底部边界行的内核已经完成，
然后才发起 halo 交换的网络通信。等待期间，\code{cudaStreamSynchronize} 会阻塞主机线程，
使其无法执行后续主机代码；同步必须由主机参与，也就意味着主机不能去完成其他有用工作。

如果通信操作本身可以放入 stream，就能摆脱这种低效：stream 机制会在通信所依赖的操作完成
后自动发起通信，无须主机线程充当同步中介。希望在 CUDA stream 的上下文中执行通信操作，
正是 NVIDIA 集体通信库（NCCL）\cite{ch23-nccl} 的设计动机之一。

NCCL 提供多种 GPU 间通信原语，既支持发送、接收等点对点通信，也支持归约、广播、散布和
聚集等集体通信。NCCL 的通信原语虽然与 MPI 相似，但有一项重要区别：NCCL 原语在 GPU
而不是主机 CPU 上运行，并且可以放入 CUDA stream。因此，使用 NCCL 后，主机线程不再需要
在内核调用与通信原语之间充当同步中介。NCCL 通信原语的另一项优势，是会感知拓扑，并针对
节点内和跨节点多 GPU 通信作高度优化。开发者因而无须为具体系统和互连自行优化通信例程。
NCCL 支持 PCIe、NVLINK、InfiniBand Verbs 和套接字等多种互连技术。

NCCL 并不能完整取代 MPI，而是可以与 MPI 结合使用、增强多 GPU 通信的库。二者结合时，
需要像建立和关闭 MPI 通信系统一样建立和关闭 NCCL 通信系统。为此，需要了解三个重要的
NCCL API：\code{ncclGetUniqueId}、\code{ncclCommInitRank} 和
\code{ncclCommDestroy}。图 \ref{fig:23-16} 给出它们的签名和作用。

\begin{figure}[htbp]
  \centering
  \begin{lstlisting}[language=C++,numbers=none,basicstyle=\ttfamily\scriptsize]
ncclResult_t ncclGetUniqueId(ncclUniqueId *uniqueId)
// 创建一个 ncclCommInitRank 所需的唯一 ID；只由一个秩调用并分发给所有秩

ncclResult_t ncclCommInitRank(ncclComm_t *comm, int nranks,
                              ncclUniqueId commId, int rank)
// 创建新通信器；comm 返回句柄，nranks 指定秩数，commId 指定唯一 ID，rank 指定调用者

ncclResult_t ncclCommDestroy(ncclComm_t comm)
// 销毁 comm 指定的通信器对象
  \end{lstlisting}
  \caption{建立和关闭通信系统所用的基本 NCCL API 函数。依据原书图 23.16 逐字符重排，
  参数说明译为中文。}
  \label{fig:23-16}
\end{figure}

图 \ref{fig:23-17} 展示把 NCCL 与 MPI 结合使用时，通常如何用图 \ref{fig:23-16} 的 API
建立和关闭 NCCL 通信系统。建立和关闭 MPI 通信系统的代码与图 \ref{fig:23-6} 相同；新增
部分包括访问 NCCL API 的头文件（第 2 行），以及建立（第 9--14 行）和关闭（第 16 行）
NCCL 通信系统的 API 调用。

\begin{figure}[htbp]
  \centering
  \begin{lstlisting}[language=C++,basicstyle=\ttfamily\scriptsize,firstnumber=1]
#include "mpi.h"
#include <nccl.h>
int main(int argc, char *argv[]) {
    MPI_Init(&argc, &argv);
    int rank;
    MPI_Comm_rank(MPI_COMM_WORLD, &rank);
    int numRanks;
    MPI_Comm_size(MPI_COMM_WORLD, &numRanks);
    ncclUniqueId ncclID;
    if(rank == 0) ncclGetUniqueId(&ncclID);
    MPI_Bcast(&ncclID, sizeof(ncclUniqueId), MPI_BYTE, 0, MPI_COMM_WORLD);
    MPI_Barrier(MPI_COMM_WORLD);
    ncclComm_t ncclComm;
    ncclCommInitRank(&ncclComm, numRanks, ncclID, rank);
    ...
    ncclCommDestroy(ncclComm);
    MPI_Finalize();
    return 0;
}
  \end{lstlisting}
  \caption{一个简单的 MPI + NCCL 主程序。依据原书图 23.17 逐字符重排。}
  \label{fig:23-17}
\end{figure}

NCCL 通信操作与 MPI 类似，同样以通信器为基础，由通信器指定参与每次操作的所有实体。因此，
一组 $n$ 个 CUDA 设备用 NCCL 相互通信以前，必须创建 NCCL 通信器。创建时，需要为通信器
中的每个 CUDA 设备分配 0 到 $n-1$ 之间的唯一秩。如果应用中的 MPI 进程或线程与 CUDA
设备一一对应，就可以直接把 MPI 秩作为 NCCL 秩，图 \ref{fig:23-17} 的主机代码正是如此。

秩分配以后，下一步是创建一个唯一对象，所有进程和线程用它同步，并确认自己属于同一个
通信器。只能由一个秩调用 \code{ncclGetUniqueId} 完成这一步。图 \ref{fig:23-17} 中，秩 0
调用该函数创建唯一 ID（第 10 行），再用 \code{MPI_Bcast} 把唯一对象从秩 0 发送给所有其他
MPI 秩（第 11 行）。注意，所有接收秩也必须执行 \code{MPI_Bcast}，才能接收广播对象。关于
\code{MPI_Bcast} 的细节，请参阅 MPI 用户指南。简而言之，它的参数依次指定被广播数据的
地址、大小和类型，执行广播的秩（本例为秩 0），以及广播范围（本例为应用的全部 MPI 秩）。
第 12 行调用 \code{MPI_Barrier}，强制所有秩相互等待，直到全部收到 ID 对象，才创建通信器。

接着，所有秩都调用 \code{ncclCommInitRank}，初始化每个秩中的 NCCL 通信器对象及通信系统
（第 14 行）。所有秩向该函数传入同一个 \code{ncclID}，从而把参与秩注册为同一个通信器的
成员；通信器在每个秩中都有一份副本。最后，主函数结束时调用 \code{ncclCommDestroy}，关闭
NCCL 通信系统（第 16 行）。本例只需一个通信器；关于在应用中使用多个 NCCL 通信器和多个
NCCL 通信系统的细节，请参阅 NCCL 用户指南\cite{ch23-nccl}。

建立好 NCCL 通信系统后，下面考察如何用 NCCL 原语替代 \code{MPI_Sendrecv} 执行 halo
交换。我们特别关心的通信原语是 \code{ncclSend} 和 \code{ncclRecv}，二者的签名与说明见图
\ref{fig:23-18}。它们与 \code{MPI_Send} 和 \code{MPI_Recv} 相似，但有一项主要区别：
二者是异步函数，并且接收 CUDA stream 对象，指定自己应在哪条 stream 中执行。

\begin{figure}[htbp]
  \centering
  \begin{lstlisting}[language=C++,numbers=none,basicstyle=\ttfamily\scriptsize]
ncclResult_t ncclSend(const void *sendbuff, size_t count,
                      ncclDataType_t datatype, int peer,
                      ncclComm_t comm, cudaStream_t stream)
// 向另一个秩发送数据；参数给出缓冲区、元素数、类型、目标秩、通信器和执行 stream

ncclResult_t ncclRecv(void *recvbuff, size_t count,
                      ncclDataType_t datatype, int peer,
                      ncclComm_t comm, cudaStream_t stream)
// 从另一个秩接收数据；参数给出缓冲区、元素数、类型、源秩、通信器和执行 stream
  \end{lstlisting}
  \caption{发送和接收数据的 NCCL API 函数。依据原书图 23.18 逐字符重排，参数说明译为中文。}
  \label{fig:23-18}
\end{figure}

第 \ref{sec:multi-gpu-mpi} 节已经说明，不能简单地用 \code{MPI_Send} 和 \code{MPI_Recv}
执行 halo 交换，因为每个秩都需要同时向一个秩发送数据、从另一个秩接收数据；为此，我们
改用把两个操作融合起来的 \code{MPI_Sendrecv}。NCCL 没有提供相应的融合发送接收函数，而是
提供了用于融合通信原语的\textbf{组调用（group call）}。NCCL 可以把一组通信原语夹在两个
API 调用之间，即 \code{ncclGroupStart} 和 \code{ncclGroupEnd}；它们的作用见图
\ref{fig:23-19}。因此，把 \code{ncclSend} 和 \code{ncclRecv} 放在两者之间，就能得到融合的
发送接收操作。NCCL 原语会在每个 CUDA 设备上并发运行，使各设备可以把边界网格点发送给
一个邻居，同时从另一个邻居接收 halo 网格点。

\begin{figure}[htbp]
  \centering
  \begin{lstlisting}[language=C++,numbers=none,basicstyle=\ttfamily\scriptsize]
ncclResult_t ncclGroupStart()
// 后续 NCCL 通信原语进入组调用，直到 ncclGroupEnd

ncclResult_t ncclGroupEnd()
// 等待自 ncclGroupStart 后的全部通信原语都已提交到指定 stream
  \end{lstlisting}
  \caption{把通信原语组织在一起的 NCCL API 函数。依据原书图 23.19 逐字符重排，
  说明文字译为中文。}
  \label{fig:23-19}
\end{figure}

现在可以用 NCCL 实现 halo 交换，并重叠计算与通信。图 \ref{fig:23-20} 给出基于 NCCL 的
主机代码。与图 \ref{fig:23-14} 的 MPI 代码相比，主要变化位于 halo 交换（第 25--36 行）。
这里，每个 \code{MPI_Sendrecv} 都替换为一组夹在 \code{ncclGroupStart} 与
\code{ncclGroupEnd} 之间的 \code{ncclSend} 和 \code{ncclRecv} 调用。更重要的是，代码不再
调用 \code{cudaStreamSynchronize()}。原先对 \code{topStream} 与 \code{bottomStream} 的
同步改为把这两条 stream 作为最后一个参数传给 \code{ncclSend} 和 \code{ncclRecv}。
CUDA 运行时保证，
只有相应 stream 中先前的内核算完 halo 行，这些通信原语才会发起。

\begin{figure}[htbp]
  \centering
  \begin{lstlisting}[language=C++,basicstyle=\ttfamily\scriptsize,firstnumber=1]
for(unsigned int iter = 0; l2norm > tol && iter < iter_max; ++iter) {

    // Reset L2 norm
    cudaMemsetAsync(l2norm_d, 0, sizeof(float), internalStream);
    cudaEventRecord(resetL2, internalStream);
    cudaStreamWaitEvent(topStream, resetL2, 0);
    cudaStreamWaitEvent(bottomStream, resetL2, 0);

    // Compute stencil
    launch_jacobi_kernel(output, input, l2norm_d, 3, nx, topStream);
    cudaEventRecord(computeTop, topStream);
    launch_jacobi_kernel(output + (ny - 3)*nx, input + (ny - 3)*nx, l2norm_d, 3,
                         nx, bottomStream);
    cudaEventRecord(computeBottom, bottomStream);
    launch_jacobi_kernel(output + nx, input + nx, l2norm_d, ny - 2, nx,
                         internalStream);

    // Copy L2 norm to host
    cudaStreamWaitEvent(internalStream, computeTop, 0);
    cudaStreamWaitEvent(internalStream, computeBottom, 0);
    cudaMemcpyAsync(l2norm_h, l2norm_d, sizeof(float), cudaMemcpyDeviceToHost,
                    internalStream);

    // Halo exchange
    const int top = rank > 0 ? rank - 1 : (numRanks - 1);
    const int bottom = (rank + 1) % numRanks;
    ncclGroupStart();
    ncclSend(output + nx, nx, ncclFloat, top, nccl_comm, topStream);
    ncclRecv(output + (ny - 1)*nx, nx, ncclFloat, bottom, nccl_comm, topStream);
    ncclGroupEnd();
    cudaEventRecord(exchangeTop, topStream);
    ncclGroupStart();
    ncclSend(output + (ny - 2)*nx, nx, ncclFloat, bottom, nccl_comm, bottomStream);
    ncclRecv(output, nx, ncclFloat, top, nccl_comm, bottomStream);
    ncclGroupEnd();
    cudaEventRecord(exchangeBottom, bottomStream);

    // Reduce L2 norm
    cudaStreamSynchronize(internalStream);
    MPI_Allreduce(l2norm_h, &l2norm, 1, MPI_FLOAT, MPI_SUM, MPI_COMM_WORLD);
    l2norm = std::sqrt(l2norm);

    cudaStreamWaitEvent(internalStream, exchangeTop, 0);
    cudaStreamWaitEvent(internalStream, exchangeBottom, 0);

    std::swap(output, input);

}
  \end{lstlisting}
  \caption{基于 NCCL、重叠计算与通信的多 GPU Jacobi 迭代法主机代码。依据原书图 23.20
  逐字符重排。}
  \label{fig:23-20}
\end{figure}

\code{MPI_Sendrecv} 与 \code{ncclSend}/\code{ncclRecv} 组调用有一项重要区别：前者是阻塞
调用，在通信完成前不会返回；后者是非阻塞调用，即使 \code{ncclGroupEnd} 已经返回，也只能
保证 \code{ncclSend} 和 \code{ncclRecv} 已经放入指定 stream。因此，必须进一步保证 halo
交换真正完成以后才开始下一次迭代。图 \ref{fig:23-20} 在 \code{topStream} 和
\code{bottomStream} 各自的 \code{ncclSend}/\code{ncclRecv} 调用后，分别记录
\code{exchangeTop} 与 \code{exchangeBottom} 事件（第 31、36 行）；进入下一次迭代以前，再
让 \code{internalStream} 等待这些事件（第 43--44 行），从而提供这项保证。

\begin{figure}[htbp]
  \centering
  \resizebox{\linewidth}{!}{%
  \begin{tikzpicture}[x=1cm,y=1cm,>=Stealth,font=\scriptsize]
    \draw[->] (0,6.4) -- (18.2,6.4) node[above left] {时间};
    \foreach \y/\name in {5.5/{主机},4.0/{顶部 stream},2.5/{底部 stream},1.0/{内部 stream}}
      \node[left,align=right] at (0,\y) {\name};
    \foreach \x/\n in {0/04,0.7/05,1.4/06,2.1/07,2.8/10,3.5/11,4.2/12,4.9/14,5.6/15,6.3/19,7/20,7.7/21,8.4/27,9.1/31,9.8/32,10.5/36}
      \node[draw,minimum width=0.7cm,minimum height=0.5cm] at (\x+0.35,5.5) {\n};
    \node[draw,minimum width=2.5cm,minimum height=0.65cm] at (14.0,5.5) {同步内部 stream};
    \node[draw,minimum width=2.2cm,minimum height=0.65cm] at (16.5,5.5) {MPI 全归约};
    \node[draw,dotted,minimum width=1.6cm] at (2.1,4.0) {等待 resetL2};
    \node[draw,minimum width=1.9cm] at (4.0,4.0) {Jacobi（顶部）};
    \node[draw,dotted,minimum width=1.0cm] at (5.55,4.0) {完成};
    \node[draw,minimum width=3.3cm] at (8.2,4.0) {NCCL 组调用：发送、接收};
    \node[draw,dotted,minimum width=1.3cm] at (11.1,4.0) {exchangeTop};
    \node[draw,dotted,minimum width=1.6cm] at (2.1,2.5) {等待 resetL2};
    \node[draw,minimum width=1.9cm] at (4.0,2.5) {Jacobi（底部）};
    \node[draw,dotted,minimum width=1.0cm] at (5.55,2.5) {完成};
    \node[draw,minimum width=3.3cm] at (8.2,2.5) {NCCL 组调用：发送、接收};
    \node[draw,dotted,minimum width=1.5cm] at (11.2,2.5) {exchangeBottom};
    \node[draw,minimum width=1.6cm] at (0.9,1.0) {memset};
    \node[draw,dotted,minimum width=1.2cm] at (2.4,1.0) {resetL2};
    \node[draw,minimum width=5.8cm] at (6.4,1.0) {Jacobi（内部）};
    \node[draw,dotted,minimum width=0.8cm] at (9.9,1.0) {等待顶};
    \node[draw,dotted,minimum width=0.8cm] at (11.1,1.0) {等待底};
    \node[draw,minimum width=1.3cm] at (12.45,1.0) {memcpy};
    \node[draw,dotted,minimum width=1.2cm] at (15.1,1.0) {等待交换顶};
    \node[draw,dotted,minimum width=1.2cm] at (17.2,1.0) {等待交换底};
    \draw[dotted,->] (2.4,1.35) -- (2.1,3.65);
    \draw[dotted,->] (2.4,1.35) -- (2.1,2.85);
    \draw[dotted,->] (11.1,3.65) -- (15.1,1.35);
    \draw[dotted,->] (11.2,2.15) -- (17.2,1.35);
    \draw[dashed,thick,->] (12.45,1.35) -- (14.0,5.15);
  \end{tikzpicture}
  }
  \caption{使用 CUDA stream 和事件重叠计算与 NCCL 通信的示例时间线。依据原书图 23.21
  重绘。}
  \label{fig:23-21}
\end{figure}

图 \ref{fig:23-21} 展示图 \ref{fig:23-20} 主机代码用 NCCL 重叠计算与通信时得到的时间线。
与图 \ref{fig:23-15} 相比，关键区别是 halo 交换不再是主机时间线中的 MPI 调用，而变成
\code{topStream} 和 \code{bottomStream} 时间线中的 NCCL 组调用；对两条 stream 的同步也
从主机时间线中的 stream 同步，变成 \code{internalStream} 时间线中的事件同步。因此，除在
\code{MPI_Allreduce} 前同步 \code{internalStream} 以外，主机已经摆脱通信和同步操作。

\section{使用 NVSHMEM 的多 GPU 模板计算}
\label{sec:multi-gpu-nvshmem}

前面已经说明如何用 MPI 和 NCCL 为 halo 交换执行数据通信。两者采用的通信风格都称为双边
通信。之所以称为“双边”，是因为一个进程向另一个进程发送数据或从中接收数据时，两个进程
都必须执行相互对应的原语。对于 MPI，发送进程的 \code{MPI_Sendrecv} 必须与接收进程的
\code{MPI_Sendrecv} 相对应；对于 NCCL，发送进程的 \code{ncclSend} 必须与接收进程的
\code{ncclRecv} 相对应。双方都要参与，是因为发送进程指定通信数据来自何处，接收进程则指定
收到的数据应放到何处。实际执行时，发送和接收进程必须等待彼此都到达这些调用以后才能发起
通信，这会产生额外延迟。

并行计算中另一种常见通信形式是\textbf{单边通信（one-sided communication）}。在单边通信
中，一个进程无须另一个进程参与，就能向它发送数据或从中接收数据。进程通过 put 操作直接
访问另一进程的地址空间，把数据存入其中；也可以通过 get 操作从另一进程的地址空间加载数据。
直接访问另一进程地址空间的一部分，无须等待对方进程就绪，可以降低数据通信延迟。单边通信
还允许在设备上执行的线程发起通信：对于设备上运行的数十万线程，配对的发送和接收调用无法
扩展，而单边 put 与 get 操作可以扩展。

OpenSHMEM\cite{ch23-openshmem} 是一种常用的单边通信库。NVIDIA 为 GPU 集群提供了
OpenSHMEM 实现，即 NVIDIA OpenSHMEM 库（NVSHMEM）\cite{ch23-nvshmem}。NVSHMEM 允许
主机线程或内核中的单个设备线程发起 put 和 get 操作。换句话说，一个进程中在 CPU 或 GPU
上运行的任何线程，都可以直接向另一个进程中另一个 GPU 的部分内存空间加载或存储数据。
在 NVSHMEM 术语中，进程或秩称为\textbf{处理元素（processing element，PE）}。

读者可能会问，发起 put 或 get 的进程中，线程如何得知目标进程地址空间里被访问数据对象的
地址。在双边通信中，这个地址由目标进程执行的配对调用指定；而在单边通信中，发起进程必须
自己指定它。为此，需要使用\textbf{对称堆（symmetric heap）}。每个进程都划出自己地址空间
的一部分，在本地容纳对称堆；所有进程都在各自对称堆本地实例的同一偏移处，分配相互对应的
内存对象。因此，发起进程要访问目标进程地址空间中的对象时，只需给出自己地址空间内对应
对象的地址。运行时可以推导对象的偏移，进而得到它在目标进程地址空间中的地址。

\begin{figure}[htbp]
  \centering
  \begin{lstlisting}[language=C++,numbers=none,basicstyle=\ttfamily\scriptsize]
void* nvshmem_malloc(size_t size)
// 在对称堆中分配至少 size 字节，返回所分配字节的地址

void nvshmem_free(void *ptr)
// 从对称堆释放 ptr 指向的数据

__device__ float nvshmem_float_g(const float* source, int pe)
// 从 PE pe 的内存中、对称地址 source 处取得一个值并返回

__device__ void nvshmem_float_p(float* dest, float val, int pe)
// 把 val 存入 PE pe 的内存中、对称地址 dest 处
  \end{lstlisting}
  \caption{为对称堆分配、释放、put 和 get 数据的 NVSHMEM API 函数。依据原书图 23.22
  逐字符重排，说明文字译为中文。}
  \label{fig:23-22}
\end{figure}

\noindent\begin{minipage}{\linewidth}
图 \ref{fig:23-22} 列出在对称堆中分配、释放以及从内核读写对象的四个 NVSHMEM API。
下面的代码用 \code{nvshmem_malloc} 对称分配保存输入、输出网格点值的数组：

\begin{lstlisting}[language=C++,numbers=none]
input = (float*) nvshmem_malloc(nx*ny*sizeof(float));
output = (float*) nvshmem_malloc(nx*ny*sizeof(float));
\end{lstlisting}
\end{minipage}

随后可以用 \code{nvshmem_free} 释放这些数组：

\begin{lstlisting}[language=C++,numbers=none]
nvshmem_free(input);
nvshmem_free(output);
\end{lstlisting}

调用这些例程时，所有 PE 都必须参与。注意，每个 PE 只为输入和输出网格分配自己的分区。
这里假定所有分区具有对称的，也就是完全相同的布局，包括顶部、底部网格点的偏移。

\begin{figure}[htbp]
  \centering
  \begin{lstlisting}[language=C++,basicstyle=\ttfamily\scriptsize,firstnumber=1]
#include <nvshmem.h>
__global__ void jacobi_kernel(float* out, float* in, float* l2norm,
                              unsigned int ny, unsigned int nx, int top, int bottom) {
    int y = blockIdx.y*blockDim.y + threadIdx.y + 1;
    int x = blockIdx.x*blockDim.x + threadIdx.x + 1;
    // Compute stencil and residue
    float residue = 0.0f;
    if (y < ny - 1 && x < nx - 1) {
        float val = 0.25*(in[y*nx + (x + 1)]
                        + in[y*nx + (x - 1)]
                        + in[(y + 1)*nx + x]
                        + in[(y - 1)*nx + x]);
        out[y*nx + x] = val;
        residue = val - in[y*nx + x];
        if (y == 1) {
            nvshmem_float_p(out + (ny - 1)*nx + x, val, top);
        }
        if (y == ny - 2) {
            nvshmem_float_p(out + x, val, bottom);
        }
    }
    // Reduce L2 norm
    float blockL2norm = blockReduce(residue*residue);
    if (threadIdx.y == 0 && threadIdx.x == 0) {
        atomicAdd(l2norm, blockL2norm);
    }
}
  \end{lstlisting}
  \caption{使用 NVSHMEM put 函数执行 halo 交换的 Jacobi 迭代模板内核。依据原书图 23.23
  逐字符重排。}
  \label{fig:23-23}
\end{figure}

图 \ref{fig:23-23} 展示如何修改图 \ref{fig:23-1} 的内核，在内核内部用 NVSHMEM 交换
halo 值。主要变化是包含 NVSHMEM 库头文件（第 1 行），把顶部和底部 PE 的 ID 作为内核参数
传入（第 3 行），以及调用 \code{nvshmem_float_p} 向顶部、底部 PE 发送边界值（第 15--20 行）。

具体而言，每个线程检查自己是否在计算本 PE 顶部边界行中的值，即检查 $y=1$（第 15 行）。
若是，线程调用 \code{nvshmem_float_p} 把该值发给顶部 PE：第二个参数传入数值，第三个参数
传入顶部 PE 的 ID（第 16 行）。由于该值需要放入顶部 PE 的底部 halo 行，线程把自己底部
halo 行中相应对称元素的地址作为第一个参数传入。类似地，每个线程还检查自己是否在计算本
PE 底部边界行中的值，即检查 $y=ny-2$（第 18 行）。若是，线程调用
\code{nvshmem_float_p} 把值发给底部 PE，第二个参数是数值，第三个参数是底部 PE 的 ID
（第 19 行）。由于该值需要放入底部 PE 的顶部 halo 行，线程把自己顶部 halo 行中相应对称
元素的地址作为第一个参数传入。

实现使用 NVSHMEM 的内核以后，下面考察调用该内核的主机代码。NVSHMEM 与 NCCL 类似，
可以和 MPI 结合使用。图 \ref{fig:23-24} 展示通常如何结合 MPI 建立 NVSHMEM。建立和关闭
MPI 通信系统的代码与图 \ref{fig:23-6} 相同；新增部分包括访问 NVSHMEM API 的头文件
（第 2--3 行），以及分配、初始化（第 10--13 行）和释放（第 15 行）NVSHMEM 库资源的 API。

\begin{figure}[htbp]
  \centering
  \begin{lstlisting}[language=C++,basicstyle=\ttfamily\scriptsize,firstnumber=1]
#include "mpi.h"
#include <nvshmem.h>
#include <nvshmemx.h>
int main(int argc, char *argv[]) {
    MPI_Init(&argc, &argv);
    int rank;
    MPI_Comm_rank(MPI_COMM_WORLD, &rank);
    int numRanks;
    MPI_Comm_size(MPI_COMM_WORLD, &numRanks);
    nvshmemx_init_attr_t attr;
    MPI_Comm mpi_comm = MPI_COMM_WORLD;
    attr.mpi_comm = &mpi_comm;
    nvshmemx_init_attr(NVSHMEMX_INIT_WITH_MPI_COMM, &attr);
    ...
    nvshmem_finalize();
    MPI_Finalize();
    return 0;
}
  \end{lstlisting}
  \caption{一个简单的 MPI + NVSHMEM 主程序。依据原书图 23.24 逐字符重排。}
  \label{fig:23-24}
\end{figure}

具体来说，代码创建 \code{nvshmemx_init_attr_t} 类型的对象（第 10 行），为初始化 NVSHMEM
通信系统提供配置细节。关于该类型支持的全部配置属性字段，请参阅 NVSHMEM 用户指南。本例
只使用属性字段 \code{mpi_comm}，它指向一个 MPI 通信器对象，指定参与 NVSHMEM 通信的进程
范围；模板示例的范围是 MPI 通信系统中的全部秩（第 11--12 行）。随后调用
\code{nvshmemx_init_attr}，用作为第二个参数传入的对象初始化 NVSHMEM 库资源（第 13 行）。
第一个参数 \code{NVSHMEMX_INIT_WITH_MPI_COMM} 表示通过 MPI 通信器初始化 NVSHMEM。最后，
主机代码完成 Jacobi 计算后，调用 \code{nvshmem_finalize} 释放 NVSHMEM 库资源（第 15 行）。
NVSHMEM 也可以在没有 MPI 通信器的情况下初始化；关于不与 MPI 结合时的初始化方法，请参阅
NVSHMEM 文档\cite{ch23-nvshmem}。

图 \ref{fig:23-25} 给出使用 NVSHMEM 内核执行 Jacobi 迭代法的主循环。复位 $L_2$ 范数平方
和（第 4 行）与归约它（第 12--16 行）的代码与前面的示例相同，主要区别在启动计算内核和
执行 halo 交换的第 7--9 行。MPI 示例图 \ref{fig:23-14} 与 NCCL 示例图 \ref{fig:23-20}
把三个内核放入三个 stream；这里则只把一个内核放入名为 \code{stream} 的一条 stream，用来
计算整个分区（第 9 行）。内核新增 \code{top} 和 \code{bottom} 两个参数，告诉它应把边界值
发送给哪些 PE。这两个值的计算方式与前面的示例相同（第 7--8 行）。本节稍后会讨论，每次
Jacobi 迭代只启动一个内核，同时执行计算和通信，有什么影响。

\begin{figure}[htbp]
  \centering
  \begin{lstlisting}[language=C++,basicstyle=\ttfamily\scriptsize,firstnumber=1]
for(unsigned int iter = 0; l2norm > tol && iter < iter_max; ++iter) {

    // Reset L2 norm
    cudaMemsetAsync(l2norm_d, 0, sizeof(float), stream);

    // Compute stencil and exchange halos
    const int top = PE > 0 ? PE - 1 : (numPEs - 1);
    const int bottom = (PE + 1) % numPEs;
    launch_jacobi_kernel(output, input, l2norm_d, ny, nx, top, bottom, stream);

    // Copy and reduce L2 norm
    cudaMemcpyAsync(l2norm_h, l2norm_d, sizeof(float), cudaMemcpyDeviceToHost,
                    stream);
    cudaStreamSynchronize(stream);
    MPI_Allreduce(l2norm_h, &l2norm, 1, MPI_FLOAT, MPI_SUM, MPI_COMM_WORLD);
    l2norm = std::sqrt(l2norm);

    nvshmemx_barrier_all_on_stream(stream);

    std::swap(output, input);

}
  \end{lstlisting}
  \caption{基于 NVSHMEM、重叠计算与通信的多 GPU Jacobi 迭代法主机代码。依据原书图 23.25
  逐字符重排。}
  \label{fig:23-25}
\end{figure}

把同时执行计算和 halo 交换的内核放入 stream 后，必须保证边界值已经到达目标进程，才能
进入下一次迭代。图 \ref{fig:23-14} 的 MPI 实现由 \code{MPI_Sendrecv} 的阻塞性质提供保证：
数据到达前调用不会返回。因此，图 \ref{fig:23-15} 时间线中的所有进程到达
\code{MPI_Allreduce} 时，可以确信全部 halo 交换数据都已收到。

图 \ref{fig:23-20} 的 NCCL 实现通过在相应 stream 中同步 \code{ncclSend} 和
\code{ncclRecv} 事件提供保证。迭代结束前的 \code{cudaStreamWaitEvent()} 调用保证每个进程
都已收到 halo 交换数据，随后才为下一次迭代复位 $L_2$ 范数平方和。

图 \ref{fig:23-25} 的 NVSHMEM 实现则由计算网格点的同一个内核发送边界值。所有已启动线程
完成网格点计算和 halo 数据发送后，内核就会结束。同步 stream 只能保证内核已经完成计算并
发出了数据，不能保证数据已经收到。为保证接收完成，需要在同一条 stream 上调用 NVSHMEM
API \code{nvshmemx_barrier_all_on_stream}（第 18 行）。这个调用让所有 PE 都在屏障处等待，
直到 stream 中任何先前内核所发起的全部 NVSHMEM 内存访问操作都已经完成。因此，所有进程
进入下一次迭代时，可以确信全部 halo 交换数据都已收到，能够安全地继续。

用 NVSHMEM 在内核中直接对 halo 值执行 put 操作有两项主要优势。第一，它自动重叠计算与
通信。采用 MPI 和 NCCL 时，必须先让内核结束才能开始通信，因此需要用不同内核分别计算
边界值与内部值：只等待边界值内核，然后发起通信，并与内部值内核并行执行。NVSHMEM 不必
启动多个内核，也不必协调它们之间的并发。部分 GPU 线程一算出 halo 值，就能在内核内部用
put 操作发送；put 发起的通信在后台进行时，其他线程继续计算内部网格点值。

第二项优势是，在内核内直接执行 put 和 get 可以把多项操作融合进一个内核。在 Jacobi 迭代法
示例中，这项优势不太明显，因为只有一种操作，也就是模板操作。然而，如果计算包含多项计算，
而且各项计算之间还会发生通信，MPI 和 NCCL 就要求用独立内核执行每项计算；NVSHMEM 则
允许把这些操作融合进同一个内核，并从内核内部通信。对于较小的内核，这种\textbf{内核融合
（kernel fusion）}可以减少内核启动开销；如果被融合的不同操作会复用数据，还可以把数据
一直保留在共享内存和寄存器中。

一个值得继续研究的主题是，图 \ref{fig:23-23} 在内核中以细粒度方式使用 put 操作。对于通过
NVLINK 连接的 GPU，这种方法性能尚可，因为 NVLINK 能提供低延迟、高带宽通信；但对于速度
更慢的互连和跨网络通信，执行更少、更粗粒度的 put 操作会更高效。NVSHMEM 提供另一函数
\code{nvshmemx_float_put_block}，同一个线程块中的线程可以协作执行一次粗粒度 put。关于这些
高级技术的用法，请参阅 NVSHMEM 文档\cite{ch23-nvshmem}。

\section{小结}
\label{sec:multi-gpu-summary}

本章介绍了使用 MPI、NCCL 和 NVSHMEM 三种编程模型，为 GPU 集群编写多 GPU 程序的基本
模式。我们以模板模式说明多个进程如何在执行期间持续交换数据。第一个实现只使用 MPI：每个
进程先计算自己分区中的所有网格点值，全部计算结束后，再把边界值发送给相邻进程。随后我们
改进 MPI 实现，先计算边界值，再把边界值发送给相邻进程，同时计算内部网格点，从而重叠
计算与通信。要实现这种重叠，需要使用 CUDA stream 并发执行多个内核。

接下来，我们考察如何把 NCCL 与 MPI 结合起来，实现同样的计算通信重叠。NCCL 的优势在于，
它允许把通信操作放入 CUDA stream，使主机无须在 GPU 内核与 MPI 通信原语之间执行同步。
此外，NCCL 库已经针对多 GPU 通信作了高度优化，程序员无须为底层系统组织和网络拓扑自行
定制通信原语。

随后，我们考察如何用 NVSHMEM 实现相同的重叠。MPI 和 NCCL 依赖双边通信，NVSHMEM 则
依赖单边通信，免除了发送操作必须与相应接收操作配对的等待延迟。此外，NVSHMEM 的 put 和
get 可以从 GPU 内核内部执行，从而允许内核融合，并透明地重叠计算与通信。

本章比较 NCCL 与 NVSHMEM 时，重点对照了 NCCL 由主机发起的双边通信和 NVSHMEM 由设备
发起的单边通信。不过，多 GPU 编程 API 仍在不断演进，彼此共有的特性也越来越多。例如，
NCCL 最近已经增加对称内存和设备发起通信的支持，将来还可能支持单边通信。另一方面，
NVSHMEM 不仅支持设备发起通信，也支持主机发起通信。通信数据粒度较粗，而且可以独立于与
通信重叠的计算来确定时，主机发起通信的效果很好；通信数据粒度较细，并且是在与通信重叠的
计算过程中才确定时，则更适合设备发起通信。

多 GPU 编程虽然与 GPU 内核编程很不相同，但本章介绍的主要编程概念，如 SPMD、秩和屏障，
在 CUDA 编程模型中都有对应概念。这印证了我们的看法：扎实掌握一种模型的并行编程以后，
学生就能快速、轻松地学习其他模型。建议读者以本章为基础，继续研究 MPI、NCCL 和 NVSHMEM
的高级功能，并把它们应用到其他重要模式中。

\section{练习}
\label{sec:multi-gpu-exercises}

\begin{enumerate}[label=\arabic*.,leftmargin=2.7em]
  \item 假设把本章的五点模板应用于一个网格，其 $x$ 维包含 64 个网格点，$y$ 维包含 512 个
  网格点；计算划分到 16 个 MPI 秩。
  \begin{enumerate}[label=\alph*.,leftmargin=2.2em]
    \item 每个进程计算多少个输出网格点值？
    \item 每个进程需要多少个 halo 网格点？
    \item 图 \ref{fig:23-12} 阶段 1 中，每个进程计算多少个边界网格点？
    \item 图 \ref{fig:23-12} 阶段 2 中，每个进程计算多少个内部网格点？
    \item 图 \ref{fig:23-12} 阶段 2 中，每个进程发送多少字节？
  \end{enumerate}

  \item 如果以下 MPI 调用传输了 4000 字节：
  \begin{center}
    \small\code{MPI_Send(ptr_a, 1000, MPI_FLOAT, 2000, 4, MPI_COMM_WORLD)}
  \end{center}
  每个被发送数据元素的大小是多少？
  \begin{enumerate}[label=\alph*.,leftmargin=2.2em]
    \item 1 字节
    \item 2 字节
    \item 4 字节
    \item 8 字节
  \end{enumerate}

  \item 下列哪些说法正确？
  \begin{enumerate}[label=\alph*.,leftmargin=2.2em]
    \item \code{MPI_Send()} 默认是阻塞的。
    \item \code{MPI_Recv()} 默认是阻塞的。
    \item MPI 消息必须至少有 128 字节。
  \end{enumerate}
  \par\noindent\textbf{译者注：}底本将本题写成单数问法“Which of the following
  statements is true?”，但 \code{MPI_Send()} 和 \code{MPI_Recv()} 默认均为阻塞调用，
  因而 a、b 同时正确；译稿相应改为复数问法。
\end{enumerate}

\begin{chapterbibliography}{9}
  \bibitem{ch23-mpi}
  W. Gropp, E. Lusk, A. Skjellum, \emph{Using MPI: Portable Parallel Programming
  with the Message-Passing Interface}, vol.~1, MIT Press, 1999.

  \bibitem{ch23-nccl}
  NVIDIA, Overview of NCCL, \url{https://docs.nvidia.com/deeplearning/nccl/user-guide/docs/overview.html},
  2024.

  \bibitem{ch23-nvshmem}
  NVIDIA, NVIDIA OpenSHMEM library (NVSHMEM) documentation,
  \url{https://docs.nvidia.com/nvshmem/api/index.html}, 2024.

  \bibitem{ch23-multigpu-models}
  NVIDIA, Multi GPU programming models, 2025.\newline
  \url{https://github.com/NVIDIA/multi-gpu-programming-models}.

  \bibitem{ch23-cub}
  NVIDIA, CUB, \url{https://nvidia.github.io/cccl/cub/index.html}, 2025.

  \bibitem{ch23-mpich}
  MPICH, MPICH, \url{https://www.mpich.org}, 2025.

  \bibitem{ch23-openmpi}
  OpenMPI, Open MPI: open source high performance computing,
  \url{https://www.open-mpi.org}, 2025.

  \bibitem{ch23-mvapich}
  MVAPICH, MVAPICH: MPI over InfiniBand, omni-path, Ethernet/iWARP, RoCE, and
  slingshot, \url{https://mvapich.cse.ohio-state.edu}, 2025.

  \bibitem{ch23-openshmem}
  B. Chapman, T. Curtis, S. Pophale, S. Poole, J. Kuehn, C. Koelbel, L. Smith,
  Introducing openshmem: shmem for the pgas community, in: \emph{Proceedings of
  the Fourth Conference on Partitioned Global Address Space Programming Model},
  2010, pp.~1--3.
\end{chapterbibliography}
